% Edit this file directly. It was generated once from paper.tex (commit ad5eb72) by build_p3.py, which is now frozen.
\newcommand{\papertitle}{Exact and perturbative criteria for finite nonnegative linearization}
\newcommand{\papershort}{EXACT AND PERTURBATIVE CRITERIA FOR FINITE LINEARIZATION}
\newcommand{\paperdate}{October 23, 2026}
\documentclass[twoside,10pt]{article}
\usepackage[T1]{fontenc}
\usepackage{lmodern}
\usepackage{amsmath,amssymb,amsthm,amsfonts,mathtools}
\usepackage[letterpaper,total={5.5in,8in},vmarginratio=1:1,hmarginratio=1:1,headsep=21pt]{geometry}
\usepackage{xcolor}
\usepackage{xr-hyper}
\usepackage[breaklinks=true,linktocpage=true]{hyperref}
\hypersetup{colorlinks=true,linkcolor=purple,urlcolor=purple,citecolor=blue!80!black,
pdfauthor={Leslie P. Polzer},pdftitle={\papertitle},
pdfsubject={Finite orthogonal polynomials and nonnegative linearization},
pdfkeywords={orthogonal polynomials, Krawtchouk polynomials, nonnegative linearization, quadratic transformation, hypergroups, Lean}}
\medmuskip=4mu
\thickmuskip=8mu
\setlength{\emergencystretch}{2em}
\numberwithin{equation}{section}
\newtheoremstyle{slbody}{8pt}{8pt}{\normalfont\slshape}{}{\bfseries}{.}{0.5em}
{\thmname{#1}\thmnumber{ #2}\thmnote{ {\normalfont\itshape(\ignorespaces#3\unskip)}}}
\theoremstyle{slbody}
\newtheorem{theorem}{Theorem}[section]
\newtheorem{lemma}[theorem]{Lemma}
\newtheorem{corollary}[theorem]{Corollary}
\newtheorem{proposition}[theorem]{Proposition}
\theoremstyle{definition}
\newtheorem{definition}[theorem]{Definition}
\theoremstyle{remark}
\newtheorem{remark}[theorem]{Remark}
\renewcommand{\qedsymbol}{$\blacksquare$}
\newcommand{\qint}[1]{[#1]_q}
\title{\papertitle}
\author{Leslie P. Polzer}
\newcommand{\authoraffiliation}{Independent Researcher}
\newcommand{\authoremail}{polzer@fastmail.com}
\date{\paperdate}
\makeatletter
\renewcommand{\section}{\@startsection{section}{1}{\z@}{-3.5ex \@plus -1ex \@minus -.2ex}{2.3ex \@plus .2ex}{\normalfont\large\bfseries}}
\renewcommand{\@seccntformat}[1]{\csname the#1\endcsname.\enspace}
\def\tagform@#1{\maketag@@@{(\ignorespaces{\oldstylenums{#1}}\unskip\@@italiccorr)}}
\renewcommand{\eqref}[1]{\textup{{\normalfont(\oldstylenums{\ref{#1}}}\normalfont)}}
\newcommand{\shorttitle}{\papershort}
\newcommand{\shortauthors}{LESLIE P. POLZER}
\newcommand{\ps@bookheader}{%
\renewcommand\@oddfoot{\hfil}\renewcommand\@evenfoot{\hfil}%
\renewcommand\@oddhead{\ifnum\value{page}>1{\small\hfil\shortauthors}\hfil\thepage\else\hfil\fi}%
\renewcommand\@evenhead{\thepage\hfil\small\shorttitle\hfil}}
\pagestyle{bookheader}
\renewcommand{\maketitle}{\begin{center}{\large\bfseries\@title}\\\vskip36pt
{\scshape\@author}\par\medskip{\small\authoraffiliation}\par\smallskip
{\small\texttt{\authoremail}}\par\smallskip{\small\@date}\par\vskip36pt\end{center}}
\makeatother
\renewenvironment{abstract}{\quotation\noindent\small{\bfseries\abstractname.\enspace}}{\endquotation}

\externaldocument[P1-]{p1}
\begin{document}
\maketitle

\begin{abstract}
We give two sufficient criteria stated directly in the coefficients of a finite stochastic three-term recurrence. The first assumes that the sum of its diagonal and downward coefficients has positive geometric increments with ratio $q\geq1$, and that the ratio of those coefficients is constant and nonnegative. A finite algebra construction gives nonnegative linearization, including products beyond the terminating degree. This includes the full positive Krawtchouk range and the dual $q$-Krawtchouk family with nonconstant increments. The underlying quadratic algebra occurs in Koelink's work on type B Hecke algebras. The second criterion gives explicit perturbation inequalities admitting a relatively open neighborhood of recurrence arrays around every Krawtchouk system in the interior of the positive linearization range. A graded version of the algebra proves positivity of Szwarc's symmetric quadratic system for $p_0\leq1/2$ and gives explicit symmetric neighbourhoods. Szwarc's inequality method survives for all products except those whose pairwise index sums exceed the size, and a reversal identity settles these deep products for palindromic systems. Near the symmetric Krawtchouk endpoint, uniform polynomial and analytic deformation rules that keep nonnegative linearization must satisfy two first-order equalities, and the stability radius of the perturbation criterion has the optimal order in the bias.
\vskip5pt
\noindent\textbf{Keywords.}\enspace Finite orthogonal polynomials; nonnegative linearization; Krawtchouk polynomials; dual $q$-Krawtchouk polynomials; quadratic transformation; polynomial hypergroups.
\vskip5pt
\noindent\textbf{MSC2020 Classification.}\enspace 33C45; 33D45; 42C05.
\end{abstract}

\vskip50pt
\baselineskip=13.5pt
\section{Introduction}\label{sec:introduction}

Nonnegative linearization (property (P)) asks whether the product of two orthogonal polynomials is a nonnegative linear combination of polynomials in the same basis. For a finite system, the identity is understood on the orthogonality support. This distinction is essential when the formal product degree exceeds the terminal degree. In a normalization with $R_i(1)=1$, nonnegative coefficients have sum one and define the product probabilities of a finite polynomial hypergroup.

The Krawtchouk case is classical. Eagleson~\cite{eagleson} studied positive definite sequences for these polynomials in 1969; Szwarc~\cite[\S5, Problem~2]{szwarc} records the nonnegative-linearization range as $0<p_0\leq1/2$ for binomial weights $\binom Njp_0^{N-j}(1-p_0)^j$ when normalized at $j=0$.\footnote{The linked author version, \S5, Problem~2, printed p.~125, uses inconsistent parameter conventions: its displayed weights have mean $N(1-p)$, whereas its degree-zero recurrence gives $K_1(x)=(x-pN)/((1-p)N)$. Thus the displayed weights give $\int K_1\,d\mu=(1-2p)/(1-p)$. Corollary~\ref{cor:kraw} derives our weights and normalization directly.} Our proof below fixes that convention directly from a two-point model. It does not treat Krawtchouk positivity itself as a new result.

Recurrence methods also have a long history. In the second note of his 1969 report, Askey~\cite[Part~II, Theorem~1 and Corollary~1]{askey} gives finite recurrence conditions leading to a hyperbolic difference-equation maximum principle. The stated polynomial conclusion concerns products with degree sum at most $N$. Szwarc's later characterization~\cite[Theorem~1]{szwarc2003} treats finite as well as infinite systems through an auxiliary discrete hyperbolic problem. His 2005 survey nevertheless asks for criteria for finite systems, including Krawtchouk and its quadratic system $Q_n$ (Section~\ref{sec:szwarc})~\cite[\S5, Problem~2]{szwarc}. Here we seek structural assumptions on the recurrence coefficients that imply positivity for every finite product.

Koelink~\cite[Theorems~3.1(ii) and 7.8]{koelink} constructed a triangular quadratic algebra associated with the type B Hecke algebra, and identified the corresponding $q$-Krawtchouk spherical functions. The dual recurrence in that construction is the geometric family used below. Our algebraic proof extracts its positive multiplication mechanism and expresses the resulting sufficient condition as two scalar patterns in a stochastic recurrence: a constant ratio and geometric increments. No priority is claimed for the algebra or for the resulting named family. For comparison, Miclo~\cite[Proposition~3.3]{miclo} gives a spectral characterization of the hypergroup property at a boundary point for irreducible finite birth--death chains by entrywise nonnegativity of products of transition matrices shifted by eigenvalues of their leading principal submatrices.

The second criterion is quantitative stability. Around each Krawtchouk array in the interior of the positive linearization range ($0<r<1$, equivalently $0<p_0<1/2$), we exhibit an explicit radius in the space of recurrence coefficients within which all finite products remain nonnegative. Its proof uses one positive term to obtain a reference margin, an error recurrence to control perturbations, and orthogonality to preserve exact zeros outside the triangle support. The resulting neighborhoods contain independently varying coefficient arrays, rather than only a reparametrized classical family. The radius is deliberately conservative and depends on $N$; no uniform stability radius at the symmetric endpoint is asserted.

Section~\ref{sec:criterion} states the geometric criterion. Sections~\ref{sec:algebra} and~\ref{sec:normalization} give its algebraic proof. Section~\ref{sec:perturbation} proves the perturbation inequalities, including explicit nongeometric examples. Section~\ref{sec:applications} fixes the Krawtchouk normalization and compares the geometric recurrence with Koelink's construction. Section~\ref{sec:szwarc} treats Szwarc's quadratic system by a graded algebra and gives explicit symmetric neighbourhoods. Section~\ref{sec:low} adapts Szwarc's inequality method to finite systems, away from the deepest products, and settles the deep products of palindromic systems. Section~\ref{sec:limits} proves that uniformly sampled polynomial deformation rules, and analytic rules whose defining functions extend holomorphically to an open neighborhood of the region $\Omega$ specified there, must satisfy two first-order equalities if they keep property (P) at the symmetric Krawtchouk systems, shows that the stability radius of Theorem~\ref{thm:perturbation} has the optimal order in $1-r$, and records counterexamples to inequality shape conditions at smaller $r$. The final section lists questions left by these results. The companion paper~\cite{P1} proves an inequality criterion that contains the geometric criterion, covers $Q_n$ through a quadratic decomposition, and shows that the stability radius decays exponentially in $N$ for every fixed $r<1$; \cite{P2} admits ratios $c_i/(b_i+c_i)$ that increase with the index.

\section{Statement and conventions}\label{sec:criterion}
For integers $j\geq0$ put $[j]_q=1+q+\cdots+q^{j-1}$, with $[0]_q=0$; in particular $[j]_1=j$.
Consider real coefficients indexed by $0\leq i\leq N$, where $N\geq1$, satisfying
\begin{equation}\label{eq:basic}
\begin{gathered}
a_0=1,\quad b_0=c_0=0,\quad a_N=0,\quad a_i+b_i+c_i=1,\\
a_i>0\ (i<N),\qquad c_i>0\ (i>0).
\end{gathered}
\end{equation}
Set $R_{-1}=0$, $R_0=1$, and determine $R_1,\ldots,R_N$ by
\begin{equation}\label{eq:recurrence}
xR_i=a_iR_{i+1}+b_iR_i+c_iR_{i-1}\qquad(0\leq i<N).
\end{equation}
At $i=0$, \eqref{eq:recurrence} gives $R_1=x$; induction at $x=1$, using the row sums and $a_i>0$ for $i<N$, gives $R_i(1)=1$ for $0\leq i\leq N$.
The terminal relation is $(x-b_N)R_N-c_NR_{N-1}=0$ on the orthogonality support. Linearization means equality of functions on this support, rather than an unrestricted polynomial identity when the degree sum exceeds $N$.

\begin{samepage}
\begin{theorem}[Geometric increment criterion]\label{thm:criterion}
Assume \eqref{eq:basic}. Set $d_i=b_i+c_i$. Suppose
\begin{enumerate}
\item $b_i/c_i=\beta$ for every $1\leq i\leq N$, where $\beta\geq0$;
\item $\delta_i=d_i-d_{i-1}>0$ for $1\leq i\leq N$;
\item for some $q\geq1$, $\delta_{i+1}=q\delta_i$ whenever $1\leq i<N$.
\end{enumerate}
Then the recurrence has a positive orthogonality measure on exactly $N+1$ distinct real points. Its normalized polynomials satisfy
\begin{equation}\label{eq:positive}
R_mR_n=\sum_{k=0}^N g(m,n;k)R_k,
\qquad g(m,n;k)\geq0,\qquad\sum_{k=0}^N g(m,n;k)=1
\end{equation}
on that support for all $0\leq m,n\leq N$.
\end{theorem}
\end{samepage}

\begin{lemma}[Coefficient parametrization]\label{lem:parameters}
The hypotheses of Theorem~\ref{thm:criterion} are equivalent to
\begin{equation}\label{eq:parameters}
a_i=\frac{q^i[N-i]_q}{[N]_q},\qquad
b_i=(1-r)\frac{[i]_q}{[N]_q},\qquad
c_i=r\frac{[i]_q}{[N]_q}
\end{equation}
for $q\geq1$, $0<r\leq1$. When $N=1$, one may choose $q=1$.
\end{lemma}
\begin{proof}
The terminal sum gives $d_N=1$, and $d_0=0$. Geometric increments give $d_i=\delta_1[i]_q$, hence $\delta_1=1/[N]_q$. The constant ratio gives $c_i=d_i/(1+\beta)$ and $b_i=\beta d_i/(1+\beta)$. Put $r=(1+\beta)^{-1}$. Finally $[N]_q-[i]_q=q^i[N-i]_q$. These calculations reverse to give the converse, including $q=1$ and $N=1$.
\end{proof}

The parameter conventions used below are
\[
\begin{array}{c|l}
q & q\geq1\quad\text{(geometric increment ratio)}\\
\beta & b_i/c_i\geq0\\
r & 1/(1+\beta)\in(0,1]\\
p & 1/r=1+\beta\geq1\\
p_0 & r/(1+r)\in(0,1/2]\\
\text{symmetric Krawtchouk case} & r=1,\quad p_0=1/2
\end{array}
\]

\section{A finite algebra with positive multiplication}\label{sec:algebra}
Fix $p\geq1$ and $q\geq1$. Starting with $A_{N+1}=\mathbb R$, form recursively the commutative real algebras
\begin{equation}\label{eq:algebra}
A_j=A_{j+1}[z_j]\big/
\left(z_j^2-\left((q-1)\sum_{l=j+1}^N z_l+p-1\right)z_j-pq^{N-j}\right).
\end{equation}
Write $A=A_1$. For $S\subseteq\{1,\ldots,N\}$ let $z_S=\prod_{j\in S}z_j$, with $z_\varnothing=1$.

\begin{lemma}[Basis and positivity]\label{lem:basis}
The elements $z_S$ form a basis of $A$, and all their multiplication coefficients are nonnegative.
\end{lemma}
\begin{proof}
Division by a monic quadratic over any commutative coefficient ring gives a unique remainder of degree at most one. Applying this successively in \eqref{eq:algebra} gives the asserted basis and dimension $2^N$.
For positivity, replace a square by
\[
z_j^2=(q-1)\sum_{l>j}z_jz_l+(p-1)z_j+pq^{N-j}.
\]
Every displayed coefficient is nonnegative. This rewriting terminates: give $z_j$ weight $2^{N-j}$ and give a monomial the sum of the weights of its factors, counted with multiplicity. Each replacement strictly decreases this nonnegative integer weight, because $l>j$ implies $2^{N-l}<2^{N-j}$. The final expression is in the unique squarefree basis, proving the claim.
\end{proof}

\begin{lemma}[Positive associative inner product]\label{lem:trace}
Let $\tau$ extract the coefficient of $z_\varnothing$. Then
\begin{equation}\label{eq:trace}
\tau(z_Sz_T)=\boldsymbol1_{S=T}\,p^{|S|}q^{\sum_{j\in S}(N-j)}.
\end{equation}
Consequently $\langle u,v\rangle=\tau(uv)$ is positive definite, and multiplication by every element of $A$ is self-adjoint.
\end{lemma}
\begin{proof}
Induct downwards through the tower. If $j$ belongs to exactly one of $S,T$, their product has the form $z_jv$ with $v\in A_{j+1}$ and has zero constant coefficient in $z_j$. If it belongs to both, the constant coefficient in $z_j$ is $pq^{N-j}z_{S\setminus\{j\}}z_{T\setminus\{j\}}$ by \eqref{eq:algebra}. If it belongs to neither, pass to $A_{j+1}$. Iteration proves \eqref{eq:trace}. All diagonal factors are positive. Associativity and commutativity give $\tau((hu)v)=\tau(u(hv))$.
\end{proof}

Put $E_k=\sum_{|S|=k}z_S$ for $0\leq k\leq N$, and set $E_{-1}=E_{N+1}=0$.

\begin{lemma}[Three-term multiplication]\label{lem:three}
For $0\leq k\leq N$,
\begin{equation}\label{eq:three}
E_1E_k=[k+1]_qE_{k+1}+(p-1)[k]_qE_k+
pq^{k-1}[N-k+1]_qE_{k-1}.
\end{equation}
The last term is interpreted as zero for $k=0$.
\end{lemma}
\begin{proof}
For $k=0$ the identity reads $E_1=E_1$. For $N=1,k=1$ it is exactly $z_1^2=(p-1)z_1+p$.
For the induction on $N$, write $F_j$ for the elementary squarefree sums in $z_2,\ldots,z_N$, with $F_j=0$ outside $0\leq j\leq N-1$. Then $E_k=F_k+z_1F_{k-1}$ and $E_1=F_1+z_1$. For $1\leq k\leq N-1$, use \eqref{eq:algebra} to obtain
\[
E_1E_k=F_1F_k+z_1F_k+qz_1F_1F_{k-1}
 +(p-1)z_1F_{k-1}+pq^{N-1}F_{k-1}.
\]
Apply the induction hypothesis to $F_1F_k$ and $F_1F_{k-1}$. The coefficients of $F_{k+1}$ and $z_1F_k$ are respectively $[k+1]_q$ and $1+q[k]_q=[k+1]_q$. The coefficients of $F_k$ and $z_1F_{k-1}$ are both $(p-1)[k]_q$, since $1+q[k-1]_q=[k]_q$. The coefficient of $F_{k-1}$ is
\[
pq^{k-1}[N-k]_q+pq^{N-1}=pq^{k-1}[N-k+1]_q,
\]
which is also the coefficient of $z_1F_{k-2}$. Regrouping gives \eqref{eq:three}.
For the endpoint $k=N$, $E_N=z_1F_{N-1}$, so the same square relation gives
\[
E_1E_N=qz_1F_1F_{N-1}+(p-1)z_1F_{N-1}+pq^{N-1}F_{N-1}.
\]
The tail endpoint identity is
\[
F_1F_{N-1}=(p-1)[N-1]_qF_{N-1}+pq^{N-2}F_{N-2}.
\]
Substitution gives $(p-1)[N]_qE_N+pq^{N-1}E_{N-1}$, as required. This separately checks the terminating boundary.
\end{proof}

\begin{proposition}[Positive subalgebra]\label{prop:subalgebra}
The space $B=\operatorname{span}\{E_0,\ldots,E_N\}$ is the subalgebra generated by $E_1$. Every product $E_mE_n$ has nonnegative coefficients in this basis.
\end{proposition}
\begin{proof}
All $E_k$ are independent, since their squarefree supports are nonempty and disjoint. Equation~\eqref{eq:three} shows that $B$ is invariant under multiplication by $E_1$. Since $[k+1]_q>0$, its interior recurrence also expresses each $E_{k+1}$ as a polynomial in $E_1$. Thus $B=\mathbb R[E_1]$, proving closure under multiplication, including products of total index greater than $N$.
Expand such a product as $\sum_k h_kE_k$. Its coefficient at each squarefree monomial $z_S$ is $h_{|S|}$. Lemma~\ref{lem:basis} makes that coefficient nonnegative. There is a subset of each cardinality $0,\ldots,N$, so every $h_k$ is nonnegative.
\end{proof}

\section{Normalization and the finite orthogonal system}\label{sec:normalization}
\begin{lemma}[Positive normalization]\label{lem:normalization}
Evaluation $z_j\mapsto pq^{N-j}$ defines an algebra character $\epsilon$ of $A$. Put $t_k=\epsilon(E_k)>0$. Then
\begin{equation}\label{eq:ratio}
t_0=1,\quad t_1=p[N]_q,\qquad
\frac{t_{k+1}}{t_k}=pq^k\frac{[N-k]_q}{[k+1]_q}\quad(0\leq k<N).
\end{equation}
\end{lemma}
\begin{proof}
For $h=N-j$, substitution in the right-hand side of \eqref{eq:algebra} gives
\[
(q-1)p^2q^h[h]_q+(p-1)pq^h+pq^h
=p^2q^h\big((q-1)[h]_q+1\big)=p^2q^{2h}.
\]
Thus all defining relations hold. The generating polynomial of the $t_k$ is
$T(u)=\prod_{j=0}^{N-1}(1+pq^ju)$, whose coefficients are strictly positive.
When $q>1$, the identity $(1+pu)T(qu)=(1+pq^Nu)T(u)$ gives
$(q^k-1)t_k=p(q^N-q^{k-1})t_{k-1}$ for $1\leq k\leq N$, which is \eqref{eq:ratio}. When $q=1$, $t_k=\binom Nk p^k$ and the same ratio follows directly.
\end{proof}

\begin{proof}[Proof of Theorem~\ref{thm:criterion}]
Use Lemma~\ref{lem:parameters} and put $p=1/r\geq1$. Construct $A,B$ as above and set $X=E_1/t_1$, $U_k=E_k/t_k$. Divide \eqref{eq:three} by $t_1t_k$ and use \eqref{eq:ratio}. The three coefficients are
\begin{align*}
\frac{[k+1]_qt_{k+1}}{p[N]_qt_k}
&=\frac{q^k[N-k]_q}{[N]_q},\\
\frac{(p-1)[k]_q}{p[N]_q}&=(1-r)\frac{[k]_q}{[N]_q},\\
\frac{pq^{k-1}[N-k+1]_qt_{k-1}}{p[N]_qt_k}
&=r\frac{[k]_q}{[N]_q}.
\end{align*}
At $k=0$ omit the third expression; at $k=N$ omit the first. Consequently $U_k=R_k(X)$ for every $k\leq N$, with the stated terminal relation.

The orthogonality support follows from the real spectral theorem in finite dimensions. On $B$ use the positive inner product from Lemma~\ref{lem:trace}. Multiplication by $X$ is self-adjoint. The vector $1=E_0$ is cyclic because the $E_k$ are polynomials of successive degrees in $X$ and form a basis. In an orthogonal eigenspace decomposition of a self-adjoint operator, the cyclic span of one vector has dimension at most the number of distinct eigenvalues. Here the cyclic span has dimension $N+1$, so there are exactly $N+1$ distinct eigenvalues, and the projection of $1$ on every eigenspace is nonzero. Give each eigenvalue the squared norm of that projection as its mass. All masses are positive, sum to $\langle1,1\rangle=1$, and yield
\[
\int f(x)h(x)\,d\mu(x)=\langle f(X),h(X)\rangle
\]
for real polynomials $f,h$, identifying an algebra element with its action on $1$. The $E_k$ are orthogonal by \eqref{eq:trace}, so the $R_k$ are orthogonal for this measure.

Proposition~\ref{prop:subalgebra} and positive $t_k$ give nonnegative product coefficients for the $U_k$. Applying the character $\epsilon$ gives their sum equal to one because $\epsilon(U_k)=1$. Polynomial identities in the algebra hold on every spectral point; conversely the $R_k$, having successive degrees at most $N$, form a basis of functions on the $N+1$ distinct points. This proves uniqueness and \eqref{eq:positive} for all products.
\end{proof}

\section{A criterion for small coefficient perturbations}\label{sec:perturbation}

Write the tridiagonal recurrence matrix with the row convention
\[
J_{i,i+1}=a_i,\qquad J_{i,i}=b_i,\qquad J_{i,i-1}=c_i,
\]
with all other entries zero and indices outside $0,\ldots,N$ omitted.

\begin{lemma}[Finite recurrence realization and triangle support]\label{lem:finite}
Every recurrence satisfying \eqref{eq:basic} has a positive orthogonality measure on exactly $N+1$ real points. If $J$ is its tridiagonal recurrence matrix, then its finite product coefficients satisfy
\begin{equation}\label{eq:matrix-product}
g(m,n;k)=[R_m(J)]_{n,k}.
\end{equation}
They vanish unless $|m-n|\leq k\leq\min(m+n,N)$, regardless of their signs inside this interval.
\end{lemma}
\begin{proof}
Choose positive $\pi_i$ with $\pi_0=1$ and $\pi_{i+1}=\pi_i a_i/c_{i+1}$. For $D=\operatorname{diag}(\sqrt{\pi_i})$, $H=DJD^{-1}$ is real symmetric. Induction in \eqref{eq:recurrence} gives the row-vector identity $e_0^\mathsf T R_i(J)=e_i^\mathsf T$, hence $e_0^\mathsf T R_i(H)=\pi_i^{-1/2}e_i^\mathsf T$. In particular $e_0$ is cyclic for $H$. The same finite spectral argument as in Theorem~\ref{thm:criterion} gives $N+1$ distinct eigenvalues and positive spectral masses at $e_0$, and
\[
\int R_iR_j\,d\mu=e_0^\mathsf TR_i(H)R_j(H)e_0
=\boldsymbol1_{i=j}\pi_i^{-1}.
\]
The terminal relation holds on this support: its polynomial $Q(x)=(x-b_N)R_N-c_NR_{N-1}$ satisfies $e_0^\mathsf TQ(J)=0$. Since $e_0$ is cyclic and $Q(J)$ commutes with $J$, $Q(J)=0$.
Multiplying a basis polynomial by $x$ acts by the corresponding row of $J$, including the terminal row. Induction on powers of $x$ therefore gives \eqref{eq:matrix-product}. Orthogonality also gives
\[
g(m,n;k)=\frac{\int R_mR_nR_k\,d\mu}{\int R_k^2\,d\mu}.
\]
If one of $m,n,k$ exceeds the sum of the other two, its polynomial is orthogonal to their product, whose degree is smaller. This proves the asserted support.
\end{proof}

\begin{lemma}[A uniform positive margin in the asymmetric case]\label{lem:margin}
For the Krawtchouk recurrence \eqref{eq:parameters} with $q=1$ and $0<r<1$, put $s=\min(r,1-r)$. Every coefficient in its triangle support obeys
\begin{equation}\label{eq:margin}
g(m,n;k)\geq (s/4)^N\quad\text{if }|m-n|\leq k\leq\min(m+n,N).
\end{equation}
\end{lemma}
\begin{proof}
Specialize the algebra above to $q=1$, and put $y_i=z_i/p$, where $p=1/r$. Then $y_i^2=(1-r)y_i+r$, and $U_m=\binom Nm^{-1}\sum_{|S|=m}y_S$. The reduction of a product $y_Sy_T$ has nonnegative coefficients, each chosen term being a product of at most $N$ factors from $\{r,1-r\}$. Its coefficient is therefore at least $s^N$.
For each admissible $k$, at least one such term has degree $k$. Assume $m\geq n$. If $k\leq m$, take $T\subseteq S$, with sizes $n,m$, and retain the linear term for exactly $k-(m-n)$ overlap coordinates. This number lies between $0$ and $n$. If $k\geq m$, choose sets with overlap $m+n-k$ and union of size $k\leq N$, and retain the linear term at every overlap coordinate.
Permutation symmetry makes the coefficient of $y_U$ depend only on $|U|=k$. Hence the coefficient in the normalized elementary-symmetric basis is at least
\[
\frac{s^N\binom Nk}{\binom Nm\binom Nn}\geq\frac{s^N}{4^N}.
\]
This uses a single positive term and a denominator bound, not an explicit formula for the product coefficients.
\end{proof}

\begin{theorem}[Explicit perturbation inequalities]\label{thm:perturbation}
Let $N\geq1$, $0<r<1$, $s=\min(r,1-r)$, and
\begin{equation}\label{eq:radius}
\eta_N(r)=\frac12\left(\frac{s}{32N}\right)^N.
\end{equation}
Suppose coefficients $\widetilde a_i,\widetilde b_i,\widetilde c_i$ satisfy \eqref{eq:basic}, with $\widetilde b_i\geq0$ for every $i$, and
\begin{equation}\label{eq:closeness}
\max_{0\leq i\leq N}\left\{
\left|\widetilde a_i-\frac{N-i}{N}\right|,
\left|\widetilde b_i-\frac{(1-r)i}{N}\right|,
\left|\widetilde c_i-\frac{ri}{N}\right|\right\}\leq\eta_N(r).
\end{equation}
Then the finite system has nonnegative linearization. Throughout its triangle support,
\[
\widetilde g(m,n;k)\geq\tfrac12(s/4)^N.
\]
The product coefficients sum to one.
\end{theorem}
\begin{proof}
Use tildes for the perturbed system and untilded letters for its $q=1$ reference at $r$. Let $\eta=\eta_N(r)$ and use the maximum absolute row-sum matrix norm. Both $J$ and $\widetilde J$ are stochastic matrices, so $\|J\|_\infty=\|\widetilde J\|_\infty=1$ and $\|\widetilde J-J\|_\infty\leq3\eta$. Put $B_m=R_m(J)$, $\widetilde B_m=\widetilde R_m(\widetilde J)$ and $e_m=\|\widetilde B_m-B_m\|_\infty$. Theorem~\ref{thm:criterion} at $q=1$ and \eqref{eq:matrix-product} give $\|B_m\|_\infty=1$ for all $m\leq N$. Also $e_0=0$ and $e_1\leq3\eta$.

Subtract the two matrix recurrences for $1\leq m<N$. Writing $\Delta_m=\widetilde B_m-B_m$ gives exactly
\begin{align*}
\widetilde a_m\Delta_{m+1}
={}&(\widetilde J-\widetilde b_m I)\Delta_m-\widetilde c_m\Delta_{m-1}\\
 &+\big(\widetilde J-J-(\widetilde b_m-b_m)I\big)B_m\\
 &-(\widetilde c_m-c_m)B_{m-1}-(\widetilde a_m-a_m)B_{m+1}.
\end{align*}
Since $\eta\leq1/(2N)$, $\widetilde a_m\geq1/(2N)$. The terms not containing a $\Delta$ have combined norm at most $6\eta$. Also $0\leq\widetilde b_m,\widetilde c_m\leq1$, so
\begin{equation}\label{eq:error-recursion}
e_{m+1}\leq4Ne_m+2Ne_{m-1}+12N\eta.
\end{equation}
Set $L=8N$. Induction gives $e_m\leq\eta L^m$ for $1\leq m\leq N$: the base case follows from $3\leq L$, and, for $m\geq1$,
\[
4N+\frac{2N}{L}+\frac{12N}{L^m}
\leq4N+\frac14+\frac32\leq8N=L.
\]
Thus every matrix entry differs from its reference by at most
\[
\eta(8N)^N=\frac12(s/4)^N.
\]
Lemma~\ref{lem:margin} gives the positive lower bound on triangle entries, while Lemma~\ref{lem:finite} makes all entries outside the triangle exactly zero. Finally $\widetilde R_m(1)=1$ follows by recurrence from the stochastic coefficient sums, so evaluating each product at $1$ gives coefficient sum one. Here $1$ belongs to the support because $\widetilde J\boldsymbol1=\boldsymbol1$ and Lemma~\ref{lem:finite} identifies that support with its spectrum. No product-positivity inequality for the perturbed system has been assumed.
\end{proof}

\begin{corollary}[Finite systems outside the geometric family]\label{cor:nonclassical}
For every $N\geq2$, take the Krawtchouk coefficients at $r=1/2$, and replace only
\[
a_1\longmapsto a_1-h,\qquad b_1\longmapsto b_1+h,
\qquad 0<h\leq\frac14(64N)^{-N}.
\]
These recurrences have nonnegative linearization and do not satisfy the constant-ratio hypothesis of Theorem~\ref{thm:criterion}. More generally, Theorem~\ref{thm:perturbation} admits a relatively open neighborhood of the asymmetric Krawtchouk array in the affine space defined by the row-sum and boundary equalities in \eqref{eq:basic}.
\end{corollary}
\begin{proof}
All initial and terminal conventions and stochastic sums are preserved. The bound on $h$ is half the radius \eqref{eq:radius} at $r=1/2$, and is smaller than $a_1=(N-1)/N$, so all required signs hold. The perturbation theorem applies. At index $1$, $b_1/c_1=1+2Nh>1$, whereas at index $N$ the ratio remains $1$; these indices differ because $N\geq2$. Thus this array lies outside the geometric family. At the asymmetric reference, all free coefficients constrained to be positive are strictly positive; the strict version of \eqref{eq:closeness} therefore contains a relatively open neighborhood in that affine space; the sign conditions define the admissible region within it.
\end{proof}

\section{Applications and scope}\label{sec:applications}
\begin{corollary}[Full positive Krawtchouk range]\label{cor:kraw}
For every $N\geq1$ and $0<p_0\leq1/2$, Krawtchouk polynomials orthogonal for binomial weights $\binom Nj p_0^{N-j}(1-p_0)^j$ have nonnegative linearization when normalized at $j=0$.
\end{corollary}
\begin{proof}
To fix the convention without relying on a named recurrence, put $r=p_0/(1-p_0)\in(0,1]$, and let $Y_1,\ldots,Y_N$ be independent variables with values $1,-r$ of probabilities $p_0,1-p_0$. Their mean is zero and $Y_i^2=(1-r)Y_i+r$. The normalized elementary symmetric sums
\[
K_i(j)=\binom Ni^{-1}\sum_{|S|=i}\prod_{l\in S}Y_l
\]
depend only on the number $j$ of negative coordinates. They are orthogonal, because the expectation of two distinct squarefree monomials is zero: some independent mean-zero variable occurs once. Their weights in $j$ are exactly the stated binomial weights. With $x=N^{-1}\sum_lY_l=1-j/(N(1-p_0))$, multiplication by $x$ gives
\[
xK_i=\frac{N-i}{N}K_{i+1}+(1-r)\frac iN K_i+r\frac iN K_{i-1}.
\]
Indeed, in multiplying the elementary sum by $\sum_lY_l$, indices outside a subset give $(i+1)$ copies of each subset of size $i+1$; the linear part of $Y_l^2$ gives $(1-r)i$ copies of each subset of size $i$; and the constant part gives $r(N-i+1)$ copies of each subset of size $i-1$. Division by the binomial normalizers gives the displayed recurrence, including $i=N$. Its leading coefficient in $x$ is positive, its degree is $i$, and its value at $j=0$ is one. These properties identify the normalized orthogonal polynomials for the binomial measure. Apply Theorem~\ref{thm:criterion} with $q=1$.
The range is sharp in this endpoint normalization: if $1/2<p_0<1$, then
$r=p_0/(1-p_0)>1$ and $g(1,1;1)=b_1=(1-r)/N<0$.
For $N=1$, this coefficient follows from the terminal relation.
\end{proof}

\begin{corollary}[The dual $q$-Krawtchouk case]\label{cor:additional}
For every $N\geq2$, $q>1$, and $0<r\leq1$, the recurrence \eqref{eq:parameters} has nonnegative linearization. In the normalization $R_0=1$, $R_1=x$, it is outside the Krawtchouk recurrence class.
\end{corollary}
\begin{proof}
Its increments are $d_i-d_{i-1}=q^{i-1}/[N]_q$, so all hypotheses hold. For a Krawtchouk recurrence in this normalization the increments of $b_i+c_i$ are constant. Here their ratio is $q>1$ and $N\geq2$, so they are not constant. An affine change of the support transporting the same endpoint-normalized basis leaves the multiplication coefficients in the basis unchanged; it cannot remove this distinction.
\end{proof}

\begin{proposition}[Relation to the existing algebraic construction]\label{prop:source}
The quadratic relations \eqref{eq:algebra} are those of Koelink~\cite[Theorem~3.1(ii), relation (3.2)]{koelink}, with $z_j=u(x^j)$ and the same parameters $p,q$. The recurrence in this paper is the normalized recurrence of his Theorem~7.1, with $r=1/p$; its polynomial-in-the-spectral-variable interpretation is the dual $q$-Krawtchouk interpretation described after his Theorem~7.8.
\end{proposition}
\begin{proof}
The source square relation is $u(x^j)^2=(q-1)u(x^j)\sum_{l>j}u(x^l)+(p-1)u(x^j)+pq^{N-j}$, which is \eqref{eq:algebra}. In Theorem~7.1 the coefficients of the next, current and previous indices are respectively $pq^i(1-q^{N-i})$, $(p-1)(1-q^i)$ and $1-q^i$. Division by $p(1-q^N)$ gives \eqref{eq:parameters} for $q>1$. For $0\leq f\leq N$, the corresponding spectral coordinate is
\[
x_f=\frac{p(1-q^{N-f})-(1-q^f)}{p(1-q^N)}
=\frac{q^{N-f}-rq^f+r-1}{q^N-1},\qquad r=1/p.
\]
Writing $\phi_f(w_d)$ for Koelink's spherical-function values, we have $R_d(x_f)=\phi_f(w_d)$ for $0\leq d\leq N$: the initial value is one and the normalized recurrences agree, so induction in $d$ identifies them.
As $q\to1$, dividing numerator and denominator by $q-1$ and using the finite geometric sums gives $x_f\to1-(1+r)f/N$.
The $q=1$ case follows directly from the construction above, without division by zero. His discussion after Theorem~7.8 explicitly interchanges polynomial degree and spectral index to obtain the dual family.
\end{proof}

The geometric criterion covers the full positive Krawtchouk range, including the symmetric endpoint. The perturbation criterion allows freely varying coefficients around each Krawtchouk array in the interior of the positive linearization range; its stated radius is conservative and decreases with $N$. Both hypotheses can be checked using a number of arithmetic comparisons proportional to $N$, without forming any general product coefficients or testing polynomial matrices. The proofs use positive algebra multiplication, a lower bound from a single positive term, and a norm estimate. They provide sufficient conditions, not a characterization of all finite systems with property (P). The geometric parameter family and underlying algebra predate Szwarc's survey; this paper makes no claim of historical priority for these constructions or criteria.

\section{Szwarc's quadratic system}\label{sec:szwarc}

Problem~2 of~\cite[\S5]{szwarc} does not stop at the Krawtchouk recurrence. With his parameter written as $p_0$, Szwarc introduces the symmetric recurrence of length $2N+1$
\begin{equation}\label{eq:szwarc-q}
yQ_{2n}=(N-n)Q_{2n+1}+nQ_{2n-1},\qquad
yQ_{2n-1}=(1-p_0)Q_{2n}+p_0Q_{2n-2},
\end{equation}
records that $Q_{2n}(y)$ is the Krawtchouk polynomial $K_n(y^2)$ of his recurrence, states that the $Q_n$ have nonnegative linearization for $p_0\leq1/2$, and asks for a general criterion that would handle this case. His Theorem~3.3 for infinite symmetric systems does not apply. In the normalization \eqref{eq:szwarc-q} its monotonicity hypotheses hold, and its hypothesis $\alpha_n\leq\gamma_n$ holds at odd indices exactly when $p_0\leq1/2$; at an even index $2n$ it reads $n\leq N-n$ and fails in the upper half. Some failure of this kind is forced in every finite system, whose terminal upward coefficient vanishes. This section proves the positivity of \eqref{eq:szwarc-q} by a graded version of the algebra of Section~\ref{sec:algebra}, and then gives an explicit coefficient criterion for symmetric finite systems near \eqref{eq:szwarc-q}. A second proof, through a quadratic decomposition of symmetric finite systems and the criterion of~\cite{P1}, is given in~\cite[Section~\ref{P1-sec:szwarc}]{P1}.

Put $M=2N$, $R_{2n}(t)=Q_{2n}(\sqrt Nt)$ and $R_{2n-1}(t)=\sqrt N\,Q_{2n-1}(\sqrt Nt)$. Substitution in \eqref{eq:szwarc-q} shows that the $R_k$ satisfy \eqref{eq:recurrence}, with $N$ replaced by $M$, for the array
\begin{equation}\label{eq:qarray}
b_k=0,\qquad
(a_{2n},c_{2n})=\Big(\frac{N-n}{N},\frac nN\Big),\qquad
(a_{2n-1},c_{2n-1})=(1-p_0,\,p_0).
\end{equation}
It satisfies \eqref{eq:basic} with $M$ in place of $N$ whenever $0<p_0<1$. The two normalizations differ by positive factors and a positive dilation, so they have the same product coefficients up to positive factors.

\begin{theorem}[Szwarc's quadratic system]\label{thm:szwarc}
Let $N\geq1$ and $0<p_0\leq1/2$. The array \eqref{eq:qarray} has a symmetric positive orthogonality measure on $2N+1$ points, and for all $0\leq m,n\leq 2N$
\[
R_mR_n=\sum_{k=0}^{2N}g(m,n;k)R_k,\qquad g(m,n;k)\geq0,\qquad\sum_k g(m,n;k)=1
\]
on that support. If $1/2<p_0<1$, then $g(2,2;2)<0$.
\end{theorem}

Throughout this section put $r=p_0/(1-p_0)$, so that $0<r\leq1$ exactly when $p_0\leq1/2$. Let $A$ be the commutative algebra generated by $y_1,\ldots,y_N$ subject to $y_j^2=(1-r)y_j+r$. This is the algebra of Section~\ref{sec:algebra} with $q=1$ after the rescaling $y_j=z_j/p$ used in Lemma~\ref{lem:margin}. Its squarefree monomials $y_S$ form a basis, and for $r\leq1$ all multiplication coefficients are nonnegative. Call an element of $A$ nonnegative if its coefficients in this basis are nonnegative.

\begin{lemma}[Edge idempotents]\label{lem:edge}
For $1\leq l\leq N$ put $f_l=(r+y_l)/(1+r)$. Then $y_lf_l=f_l$ and $f_l^2=f_l$. The ideal $f_lA$ has basis $f_ly_S$, $S\subseteq\{1,\ldots,N\}\setminus\{l\}$, and $f_l u$ has nonnegative coefficients in this basis whenever $u\in A$ is nonnegative.
\end{lemma}
\begin{proof}
First $(1+r)y_lf_l=ry_l+(1-r)y_l+r=(1+r)f_l$. Then $(1+r)f_l^2=rf_l+y_lf_l=(1+r)f_l$. Let $A'$ be generated by the $y_j$ with $j\neq l$, so that $A=A'\oplus y_lA'$. Since $f_ly_l=f_l$, $f_lA=f_lA'$. The map $u\mapsto f_lu=(ru+y_lu)/(1+r)$ is injective on $A'$ by the direct sum, which gives the basis. For the last claim, $f_ly_S=f_ly_{S\setminus\{l\}}$ for every $S$.
\end{proof}

Let $V=A\oplus\bigoplus_{l=1}^Nf_lA$, on which $A$ acts by multiplication in every summand, and call a vector nonnegative if all of its coordinates are nonnegative in the bases above. Define the $A$-linear operator
\[
Y(u;v_1,\ldots,v_N)=\Big(\frac1N\sum_{l=1}^Nv_l\,;\,f_1u,\ldots,f_Nu\Big).
\]
Put $U_{2n}=\binom Nn^{-1}\sum_{|S|=n}y_S\in A$ for $0\leq n\leq N$. For $0\leq n\leq N-1$, let $U_{2n+1}$ be the vector in $\bigoplus_lf_lA$ whose $l$th coordinate is $\binom{N-1}n^{-1}\sum f_ly_S$, summed over $S\subseteq\{1,\ldots,N\}\setminus\{l\}$ with $|S|=n$.

\begin{lemma}[Three-term action]\label{lem:qthree}
With $U_{-1}=U_{2N+1}=0$ and the coefficients \eqref{eq:qarray}, $YU_k=a_kU_{k+1}+c_kU_{k-1}$ for $0\leq k\leq2N$.
\end{lemma}
\begin{proof}
The $l$th coordinate of $Y\sum_{|S|=n}y_S$ is $\sum_{S\not\ni l}f_ly_S+\sum_{S'\not\ni l}f_ly_{S'}$, with $|S|=n$ and $|S'|=n-1$, because $f_ly_S=f_ly_{S\setminus\{l\}}$. Dividing by $\binom Nn$ and using $\binom{N-1}n/\binom Nn=(N-n)/N$ and $\binom{N-1}{n-1}/\binom Nn=n/N$ gives the even case. For the odd case, $(1+r)f_ly_S=ry_S+y_{S\cup\{l\}}$ when $l\notin S$. Each $n$-set occurs for $N-n$ choices of $l$ and each $(n+1)$-set for $n+1$ choices, so
\[
YU_{2n+1}=\frac{r(N-n)\binom Nn U_{2n}+(n+1)\binom N{n+1}U_{2n+2}}{N(1+r)\binom{N-1}n}
=\frac{rU_{2n}+U_{2n+2}}{1+r}.
\]
Finally $r/(1+r)=p_0$ and $1/(1+r)=1-p_0$.
\end{proof}

\begin{proof}[Proof of Theorem~\ref{thm:szwarc}]
Lemma~\ref{lem:finite} gives a positive measure on the $2N+1$ zeros of the terminal polynomial $T(t)=tR_{2N}(t)-R_{2N-1}(t)$. Since $b\equiv0$, the matrix $J$ is conjugate to $-J$ by $\operatorname{diag}((-1)^i)$, which fixes $e_0$; hence the measure is symmetric and $R_k(-t)=(-1)^kR_k(t)$. Put $\Phi(f)=f(Y)U_0$, where $U_0=(1;0,\ldots,0)$. Lemma~\ref{lem:qthree} gives $\Phi(R_k)=U_k$ for $k\leq2N$ and $\Phi(T)=YU_{2N}-U_{2N-1}=0$. Every polynomial vanishing on the support is a multiple of $T$, so $\Phi$ maps it to zero. Hence the product identity on the support becomes
\[
R_m(Y)R_n(Y)U_0=\sum_kg(m,n;k)U_k.
\]
The supports of the $U_k$ in the bases of $V$ are nonempty and disjoint: $U_{2n}$ uses the $y_S$ with $|S|=n$, and $U_{2n+1}$ uses the $f_ly_S$ with $l\notin S$, $|S|=n$. Each basis vector lies in exactly one of them, with positive coefficient. Consequently $g(m,n;k)\geq0$ for every $k$ whenever the left side is nonnegative.

Suppose $m$ is even. Since polynomials in $Y$ commute and $Y$ is $A$-linear,
\[
R_m(Y)R_n(Y)U_0=R_n(Y)(U_mU_0)=U_m\,R_n(Y)U_0=U_mU_n.
\]
Here $U_m\in A$ and $U_n$ are nonnegative, so their product is nonnegative by the positivity of $A$ and Lemma~\ref{lem:edge}. The same holds when $n$ is even, since $g(m,n;k)=g(n,m;k)$. If $m,n$ are odd and $k$ is odd, then $g(m,n;k)=0$ by symmetry of the measure. If $m,n$ are odd and $k$ is even, then $g(m,n;k)\int R_k^2\,d\mu=\int R_mR_nR_k\,d\mu=g(k,m;n)\int R_n^2\,d\mu\geq0$. The point $t=1$ belongs to the support because $J\boldsymbol1=\boldsymbol1$, and $R_k(1)=1$, so the coefficients sum to one.

For $1/2<p_0<1$ we have $r>1$. The identities above are purely algebraic and remain valid, but $A$ is no longer positive. In particular $U_2^2=N^{-2}\big((1-r)\sum_jy_j+Nr+2\sum_{i<j}y_iy_j\big)$, whose $U_2$ coefficient $g(2,2;2)=(1-r)/N$ is negative.
\end{proof}

Since $Y^2$ acts on $A$ by multiplication by $N^{-1}\sum_lf_l=(r+X)/(1+r)$, with $X=N^{-1}\sum_ly_l$, the even part of the theorem is Corollary~\ref{cor:kraw}: $R_{2n}(t)=K_n\big((1+r)t^2-r\big)$ in the normalization of that corollary. Thus Theorem~\ref{thm:szwarc} recovers the full positive Krawtchouk range by Szwarc's quadratic route. At $p_0=1/2$ the operator $Y$ is the normalized adjacency operator of the subdivided $N$-cube seen from a vertex, and the theorem reduces to counting in that graph.

\begin{lemma}[Margin for the quadratic system]\label{lem:qmargin}
Let $N\geq1$, $0<p_0<1/2$, $r=p_0/(1-p_0)$ and $s=\min(r,1-r)$. For the array \eqref{eq:qarray},
\[
g(m,n;k)\geq(rs/8)^N\qquad\text{whenever }|m-n|\leq k\leq m+n\text{ and }m+n+k\text{ is even}.
\]
All other coefficients vanish.
\end{lemma}
\begin{proof}
The vanishing statements are Lemma~\ref{lem:finite} and the symmetry used above. We use $\binom Nk\leq2^N$ throughout. Up to the order of $m,n$, there are three cases.

If $m,n,k$ are even, the coefficient is a Krawtchouk coefficient computed in $A$, so Lemma~\ref{lem:margin} gives the lower bound $(s/4)^N\geq(rs/8)^N$.

Let $m=2a$, $n=2b+1$ and $k=2c+1$, so that
\[
\max(a-b-1,\,b-a)\leq c\leq\min(a+b,\,N-1).
\]
Fix $l$ and a set $U\not\ni l$ with $|U|=c$. The $l$th coordinate of $U_mU_n$ is
\[
\binom Na^{-1}\binom{N-1}b^{-1}\sum_{|T|=a}\ \sum_{S\not\ni l,\ |S|=b}f_l\,y_{T\setminus\{l\}}\,y_S .
\] Every nonzero term arising from a chosen reduction of $y_{T\setminus\{l\}}y_S$ is a product of at most $N-1$ factors from $\{r,1-r\}$, and therefore has coefficient at least $s^{N-1}$. Put $a'=|T\setminus\{l\}|$, which can equal $a-1$ when $a\geq1$ and $a$ when $a\leq N-1$. In the $N-1$ coordinates other than $l$, sets of sizes $a'$ and $b$ with overlap $o$, retaining the linear term on $t\leq o$ overlap coordinates, give monomials of size $a'+b-2o+t$. As $o$ ranges over its admissible values these sizes fill the interval $\big[|a'-b|,\min(a'+b,N-1)\big]$. The two admissible values of $a'$ cover every $c$ above. Hence $f_ly_U$ has coefficient at least $s^{N-1}\binom Na^{-1}\binom{N-1}b^{-1}$, and comparison with $U_k$ gives
\[
g(m,n;k)\geq\binom{N-1}c\frac{s^{N-1}}{\binom Na\binom{N-1}b}\geq\frac{2s^{N-1}}{4^N}\geq\Big(\frac s4\Big)^N.
\]

Finally let $m=2b+1$, $n=2c+1$ and $k=2a$. With $\pi_i$ as in the proof of Lemma~\ref{lem:finite}, \eqref{eq:qarray} gives $\pi_{2a}=\binom Nar^{-a}$ and $\pi_{2c+1}=\binom{N-1}cr^{-c}/p_0$. Orthogonality gives
\[
g(2b+1,2c+1;2a)=g(2a,2b+1;2c+1)\frac{\pi_{2a}}{\pi_{2c+1}}
\geq\frac{2s^{N-1}}{4^N}\cdot\frac{p_0r^N}{2^{N-1}}
\geq\frac{2s^{N-1}r^{N+1}}{8^N},
\]
using $r^{c-a}\geq r^N$ and $p_0\geq r/2$. Since $s\leq r$, the last quantity is at least $(rs/8)^N$.
\end{proof}

The proof of Theorem~\ref{thm:perturbation} uses only three properties of its reference: nonnegative linearization with sums one, a lower bound for the upward coefficients, and a margin on the known support. We record it in that generality.

\begin{lemma}[Stability of a reference with a margin]\label{lem:stability}
Let $a,b,c$ satisfy \eqref{eq:basic} on the indices $0,\ldots,M$, have nonnegative linearization, and satisfy $a_i\geq\alpha$ for $i<M$, where $0<\alpha\leq1$. Suppose $g(m,n;k)\geq\delta>0$ on a set $\Sigma$ of index triples. Let $\widetilde a,\widetilde b,\widetilde c$ satisfy \eqref{eq:basic}, with $\widetilde b_i\geq0$, and differ from $a,b,c$ by at most $\eta$ in every entry. If $\eta\leq\alpha/2$ and $\eta(8/\alpha)^M\leq\delta/2$, then $|\widetilde g(m,n;k)-g(m,n;k)|\leq\delta/2$ for all $m,n,k$; in particular $\widetilde g\geq\delta/2$ on $\Sigma$.
\end{lemma}
\begin{proof}
Repeat the proof of Theorem~\ref{thm:perturbation} with $N$ replaced by $M$. Now $\widetilde a_m\geq\alpha/2$, so the subtracted recurrences give $e_{m+1}\leq(4e_m+2e_{m-1}+12\eta)/\alpha$. With $L=8/\alpha\geq8$, induction gives $e_m\leq\eta L^m$: the base case is $e_1\leq3\eta$, and the step needs $4L^m+2L^{m-1}+12\leq8L^m$, which holds for $m\geq1$. Every entry of $\widetilde B_m-B_m$ is therefore at most $\eta L^M\leq\delta/2$, and \eqref{eq:matrix-product} identifies these entries with the product coefficients.
\end{proof}

Theorem~\ref{thm:perturbation} is the case $M=N$, $\alpha=1/N$ and $\delta=(s/4)^N$ of this lemma.

\begin{theorem}[Symmetric criterion near Szwarc's system]\label{thm:symmetric}
Let $N\geq1$, $0<p_0<1/2$, $r=p_0/(1-p_0)$, $s=\min(r,1-r)$, and
\[
\eta=\frac12\Big(\frac{rs}8\Big)^N(16N)^{-2N}.
\]
Let $c_k$ be the downward coefficients in \eqref{eq:qarray}. Take real numbers $\widetilde c_1,\ldots,\widetilde c_{2N-1}$ with $|\widetilde c_k-c_k|\leq\eta$, and put $\widetilde c_0=0$, $\widetilde c_{2N}=1$, $\widetilde b_k=0$ and $\widetilde a_k=1-\widetilde c_k$. The resulting symmetric finite system has nonnegative linearization with sums one. Moreover $\widetilde g(m,n;k)\geq\frac12(rs/8)^N$ whenever $|m-n|\leq k\leq m+n$ and $m+n+k$ is even, and $\widetilde g(m,n;k)=0$ otherwise.
\end{theorem}
\begin{proof}
The perturbed array satisfies \eqref{eq:basic} with $M=2N$: $\widetilde a_k\geq1/N-\eta>0$ for even $k<2N$, $\widetilde a_k\geq1/2-\eta>0$ for odd $k$, and $\widetilde c_k\geq\min(p_0,1/N)-\eta>0$ for $k>0$. Theorem~\ref{thm:szwarc} shows that the reference has nonnegative linearization, and $a_k\geq1/(2N)$ for $k<2N$. Lemma~\ref{lem:qmargin} gives the margin $\delta=(rs/8)^N$ on the stated set. Apply Lemma~\ref{lem:stability} with $\alpha=1/(2N)$: then $\eta\leq\alpha/2$ and $\eta(16N)^{2N}=\delta/2$. Off the stated set, $\widetilde g$ vanishes by Lemma~\ref{lem:finite} and by the symmetry of the perturbed measure, which holds because $\widetilde b\equiv0$. The sums equal one because $\widetilde R_m(1)=1$ and $1$ lies in the support.
\end{proof}

\begin{corollary}[Symmetric systems outside both families]\label{cor:symmetric-new}
For $N\geq2$, $0<p_0<1/2$ and $0<h\leq\eta$, the array obtained from \eqref{eq:qarray} by $c_1\mapsto p_0+h$, $a_1\mapsto1-p_0-h$ has nonnegative linearization. It is neither of the form \eqref{eq:qarray} nor of the form \eqref{eq:parameters}. More generally, Theorem~\ref{thm:symmetric} admits a relatively open set of symmetric stochastic arrays, with $2N-1$ free coefficients, around each array \eqref{eq:qarray} with $0<p_0<1/2$.
\end{corollary}
\begin{proof}
Theorem~\ref{thm:symmetric} applies. The odd downward coefficients $c_1=p_0+h$ and $c_3=p_0$ differ, so the array is not of the form \eqref{eq:qarray}; here $N\geq2$ ensures $3<2N$. In \eqref{eq:parameters}, on $2N+1$ indices, the downward coefficients $c_i=r[i]_q/[2N]_q$ are strictly increasing, whereas $c_3<c_1$.
\end{proof}

The quadratic decomposition~\cite[Theorem~\ref{P1-thm:decomp}]{P1} also transfers the perturbation criteria: a symmetric array has nonnegative linearization whenever its two normalized components satisfy the hypotheses of Theorem~\ref{thm:perturbation}, or those of Lemma~\ref{lem:stability} around a reference with a margin, provided the set $\Sigma$ of that lemma contains every coefficient that can be nonzero for the perturbed components, or the remaining coefficients are separately known to stay nonnegative. Lemma~\ref{lem:stability} controls only $\Sigma$: for example, starting from $K_{2,1}$, keeping $b\equiv0$ and $c_2=1$, and setting $c_1=\tfrac12+\varepsilon$, $a_1=\tfrac12-\varepsilon$ with $0<\varepsilon<\tfrac12$ gives $g(2,2;2)=-2\varepsilon/(\tfrac12-\varepsilon)<0$, although the coefficients positive at the reference remain positive. Conversely, the even component of every array admitted by Theorem~\ref{thm:symmetric} is a nonsymmetric finite system with nonnegative linearization.

\section{Inequality criteria away from the deep products}\label{sec:low}
Szwarc's inequality method~\cite[Theorems~3.2 and~3.4]{szwarc} survives the passage to finite systems for all products whose two factors have total degree at most $N$. Only the terminal coefficients are excluded from its hypotheses. Write a finite system in a positive renormalization $P_n=R_n/\sigma_n$ $(\sigma_0=1)$ as $xP_n=\gamma_nP_{n+1}+\beta_nP_n+\alpha_nP_{n-1}$, so that $\gamma_n=a_n\sigma_{n+1}/\sigma_n$, $\beta_n=b_n$, $\alpha_n=c_n\sigma_{n-1}/\sigma_n$ and $\gamma_N=0$.

\begin{theorem}[Low products]\label{thm:low}
Let $a,b,c$ satisfy \eqref{eq:basic}. Suppose that for some positive $\sigma$,
\begin{enumerate}
\item $\beta_n$ is nondecreasing for $0\leq n\leq N-1$, and $\alpha_n$ is nondecreasing for $0\leq n\leq N$;
\item $\alpha_n+\gamma_n$ is nondecreasing for $0\leq n\leq N-2$;
\item $\gamma_i\geq\alpha_j$ whenever $0\leq j\leq i$ and $i+j\leq N-2$.
\end{enumerate}
Then $g(n,m;k)\geq0$ whenever $n+m\leq N$. Consequently a product coefficient can be negative only for a triple $(n,m,k)$ in which each two indices have sum greater than $N$.
\end{theorem}
\begin{proof}
Coefficients in two positive normalizations differ by positive factors, so we may work with the polynomials $\widetilde P_n$ of~\cite[(3.9)]{szwarc}, which satisfy $x\widetilde P_n=\alpha_{n+1}\widetilde P_{n+1}+\beta_n\widetilde P_n+\gamma_{n-1}\widetilde P_{n-1}$ with $\alpha_{N+1}=0$. Fix $k$ and let $u(n,m)$ be the coefficient of $\widetilde P_k$ in $\widetilde P_n\widetilde P_m$. Multiplying the product by $x$ in the two ways shows $Hu=0$ at every point of $\{0,\ldots,N\}^2$, where $H=L_1-L_2$ is the operator of~\cite[\S3.4]{szwarc} built from these coefficients, and $u(n,0)=\boldsymbol1_{n=k}$, $u(n,m)=u(m,n)$.

We prove $u(n,m)\geq0$ for $n\geq m$, $n+m\leq N$, by induction on $m$, following~\cite[Theorem~3.4]{szwarc}. Let $v$ be the indicator of the points $(i,j)$ with $j<m$, $|i-n|<m-j$ and $n+m-i-j$ odd. Every point in the support of $v$ satisfies $j<i$ and $i+j\leq n+m-1\leq N-1$, so $\langle Hu,v\rangle=0$ involves only rows where the finite recurrence holds. Hence $0=\langle u,H^*v\rangle$, and every point in the support of $H^*v$ has $i+j\leq n+m\leq N$. With $\alpha_0=0$ and $v$ set to zero outside its support, transposition gives
\[
(H^*v)(i,j)=\alpha_iv(i-1,j)+\gamma_iv(i+1,j)+(\beta_i-\beta_j)v(i,j)-\alpha_jv(i,j-1)-\gamma_jv(i,j+1),
\]
as in~\cite[\S3.4]{szwarc}. At $(n,m)$ it is $-\alpha_m<0$. At the other points it is one of
\[
\beta_i-\beta_j,\quad (\alpha_i+\gamma_i)-(\alpha_j+\gamma_j),\quad \alpha_i-\alpha_j,\quad \gamma_i-\alpha_j,
\]
with $j\leq i$. The first occurs only for $i\leq N-1$. The second occurs only at interior points, where $i+j\leq n+m-2\leq N-2$. The third occurs on the edge $i+j=n+m$. The last occurs on the edge $i-j=n-m$, where $i+j\leq N-2$. Hypotheses (i)--(iii) make all of these nonnegative, and all such points have second coordinate less than $m$. By the induction hypothesis and $u(i,0)\geq0$, $\alpha_mu(n,m)=\sum_{(i,j)\neq(n,m)}(H^*v)(i,j)u(i,j)\geq0$. For the last statement, $g(n,m;k)$, $g(n,k;m)$ and $g(m,k;n)$ are positive multiples of $\int R_nR_mR_k\,d\mu$.
\end{proof}

\begin{corollary}[Convex profiles]\label{cor:convex-low}
Let $b_i=(1-r)d_i$, $c_i=rd_i$, $a_i=1-d_i$ with $0<r\leq1$ and a convex profile $0=d_0<d_1<\cdots<d_N=1$, that is, with nondecreasing increments. Then $g(n,m;k)\geq0$ whenever two of $n,m,k$ have sum at most $N$. This includes every array \eqref{eq:parameters}.
\end{corollary}
\begin{proof}
Take $\sigma_n=r^{n/2}$. Then $\alpha_n=\sqrt r\,d_n$ and $\gamma_n=\sqrt r\,(1-d_n)$, so $\alpha$ and $\beta=(1-r)d$ are nondecreasing and $\alpha_n+\gamma_n=\sqrt r$ is constant. Convexity and $d_0=0$ give $d_i+d_j\leq d_{i+j}$, so $\gamma_i-\alpha_j=\sqrt r\,(1-d_i-d_j)\geq\sqrt r\,(1-d_{i+j})>0$ for $i+j\leq N-2$. Theorem~\ref{thm:low} applies.
\end{proof}

For such profiles, positivity can therefore fail only on triples with all pairwise sums above $N$. The convex counterexamples recorded in Section~\ref{sec:open} fail at $(N,N,N)$, which is such a triple.

The cone argument extends to all products if the monotonicity conditions in (i)--(ii) are imposed through the terminal index $N$ and the edge inequalities in (iii) are imposed for every $0\leq j\leq i\leq N-1$.

\begin{theorem}[All products]\label{thm:all}
Let $a,b,c$ satisfy \eqref{eq:basic}, and put $\gamma_N=0$. Suppose that for some positive $\sigma$,
\begin{enumerate}
\item $\beta_n$ and $\alpha_n$ are nondecreasing for $0\leq n\leq N$;
\item $\alpha_n+\gamma_n$ is nondecreasing for $0\leq n\leq N$;
\item $\gamma_i\geq\alpha_j$ whenever $0\leq j\leq i\leq N-1$.
\end{enumerate}
Then the system has nonnegative linearization.
\end{theorem}
\begin{proof}
Use the notation of the proof of Theorem~\ref{thm:low}. Since $\alpha_{N+1}=0$, the relation $x\widetilde P_N=\beta_N\widetilde P_N+\gamma_{N-1}\widetilde P_{N-1}$ is the terminal relation, which holds on the support. So $Hu=0$ at every point of $\{0,\ldots,N\}^2$, as noted there. We prove $u(n,m)\geq0$ for all $n\geq m$ by induction on $m$. Let $v$ be the indicator of the points $(i,j)$ with $0\leq i\leq N$, $j<m$, $|i-n|<m-j$ and $n+m-i-j$ odd. Then $0=\langle Hu,v\rangle=\langle u,H^*v\rangle$, and $H^*v$ is given by the formula in that proof. At $i=N$ the term $\gamma_Nv(N+1,j)$ is absent; this is the convention $\gamma_N=0$. Hence $H^*v(n,m)=-\alpha_m$, and at every other point $(i,j)$ with $H^*v(i,j)\neq0$ we have $j<m$, $j\leq i\leq N$, and $H^*v(i,j)$ is one of
\[
\beta_i-\beta_j,\quad (\alpha_i+\gamma_i)-(\alpha_j+\gamma_j),\quad \alpha_i-\alpha_j,\quad \gamma_i-\alpha_j.
\]
The last one occurs only on the edge $i-j=n-m$, where $i\leq n-1\leq N-1$. Hypotheses (i)--(iii) make all of these nonnegative. The induction hypothesis then gives $\alpha_mu(n,m)\geq0$, and the case $n<m$ follows by symmetry.
\end{proof}

\begin{corollary}[Heavy top step]\label{cor:heavy}
Let $b_i=(1-r)d_i$, $c_i=rd_i$, $a_i=1-d_i$ with $0<r\leq1$ and an increasing profile $0=d_0<d_1<\cdots<d_N=1$ with $d_{N-1}\leq\tfrac12$. Then the system has nonnegative linearization.
\end{corollary}
\begin{proof}
Take $\sigma_n=r^{n/2}$. Then $\alpha_n=\sqrt r\,d_n$ and $\gamma_n=\sqrt r\,(1-d_n)$ for $n<N$, $\gamma_N=0$, and $\beta_n=(1-r)d_n$. So $\alpha$ and $\beta$ are nondecreasing, and $\alpha_n+\gamma_n=\sqrt r$ for every $0\leq n\leq N$, because $d_N=1$. For $j\leq i\leq N-1$, $\gamma_i-\alpha_j=\sqrt r\,(1-d_i-d_j)\geq\sqrt r\,(1-2d_{N-1})\geq0$. Theorem~\ref{thm:all} applies.
\end{proof}

The profile need not be convex. The corollary contains the arrays \eqref{eq:parameters} with $q\geq2$, since then $d_{N-1}=(q^{N-1}-1)/(q^N-1)\leq\tfrac12$, but not the Krawtchouk arrays with $N\geq3$. The bound $d_{N-1}\leq\tfrac12$ cannot be relaxed in general: for $N=2$ and $r=1$, $g(2,2;2)=(1-2d_1)/(1-d_1)$. For constant-ratio arrays with a convex profile, $d_i+d_j\leq d_{i+j}\leq1$ whenever $i+j\leq N$, so the only terms that can be negative in the proof are the edge terms $\gamma_i-\alpha_j=\sqrt r\,(1-d_i-d_j)$ with $i+j>N$. For the Krawtchouk arrays these edge coefficients are $\sqrt r\,(N-i-j)/N$; any negative weighted contributions they make to the cone identity must be compensated by its other terms.

\paragraph{Deep triples and reversal.}
Call $(n,m,k)$ deep if each two of its indices have sum greater than $N$. The reversed array $a'_i=c_{N-i}$, $b'_i=b_{N-i}$, $c'_i=a_{N-i}$ need not satisfy $a'_0=1$, $b'_0=0$. Its polynomials $R'_i$ are defined by the reversed recurrence from $R'_0=1$, so $R'_1=(x-b_N)/c_N$, and they still satisfy $R'_i(1)=1$; its matrix is $\Pi J\Pi$, where $\Pi$ reverses the order of the indices. Write $g'$ for its product coefficients.

\begin{lemma}[Reversal identity]\label{lem:reversal}
On the support, $R_NR'_t=R_{N-t}$ for $0\leq t\leq N$.
\end{lemma}
\begin{proof}
From $e_0^{\mathsf T}R_i(J)=e_i^{\mathsf T}$ for the array and for its reversal, $e_N^{\mathsf T}R'_t(J)=e_{N-t}^{\mathsf T}$. Hence $e_0^{\mathsf T}R_N(J)R'_t(J)=e_N^{\mathsf T}R'_t(J)=e_0^{\mathsf T}R_{N-t}(J)$. Both matrices commute with $J$ and $e_0$ is cyclic, so $R_N(J)R'_t(J)=R_{N-t}(J)$, which is the claim on the spectrum.
\end{proof}

\begin{theorem}[Deep reduction]\label{thm:deep}
For every array satisfying \eqref{eq:basic} and every deep triple $(x,y,z)$,
\[
g(x,y;z)=\sum_{t=0}^Ng'(N-x,N-y;t)\,g(N,N-t;z).
\]
Moreover $g(N,N;N-l)$ is the coefficient of $R'_l$ in $R_N$. Hence all coefficients are nonnegative if (a) $g\geq0$ on non-deep triples, (b) $g'(i,j;t)\geq0$ whenever $i+j<N$, and (c) $R_N$ is a nonnegative combination of $R'_0,\ldots,R'_N$.
\end{theorem}
\begin{proof}
By Lemma~\ref{lem:reversal}, $R_x=R_NR'_{N-x}$ and $R_y=R_NR'_{N-y}$ on the support. Expanding $R'_{N-x}R'_{N-y}=\sum_tg'(N-x,N-y;t)R'_t$ and using Lemma~\ref{lem:reversal} again gives $R_xR_y=\sum_tg'(N-x,N-y;t)R_NR_{N-t}$. Expanding $R_NR_{N-t}$ gives the formula. Writing $R_N=\sum_le_lR'_l$, Lemma~\ref{lem:reversal} gives $R_N^2=\sum_le_lR_{N-l}$, so $e_l=g(N,N;N-l)$. For the last statement, $(N-x)+(N-y)<N$, so (b) makes the weights $g'$ nonnegative. It remains to see $g(N,n;z)\geq0$. If $(N,n,z)$ is not deep this is (a). Otherwise $n+z>N$, and $g(N,n;z)$ is a positive multiple of $g(n,z;N)=\sum_tg'(N-n,N-z;t)\,g(N,N-t;N)$; here $g(N,N-t;N)$ is a positive multiple of $g(N,N;N-t)$, which is nonnegative by (c).
\end{proof}

\begin{theorem}[Palindromic criterion]\label{thm:palindromic}
Let $a,b,c$ satisfy \eqref{eq:basic} with $a_n=c_{N-n}$ and $b_n=b_{N-n}$ for all $n$, and satisfy the hypotheses of Theorem~\ref{thm:low} for some positive $\sigma$. Then the system has nonnegative linearization.
\end{theorem}
\begin{proof}
The reversed array equals the array, so $R'=R$ and $g'=g$. Theorem~\ref{thm:low} gives $g\geq0$ on all non-deep triples, in particular $g'(i,j;t)\geq0$ for $i+j<N$. Since $R_N=R'_N$ is a single reversed polynomial, $g(N,N;\cdot)\geq0$. Theorem~\ref{thm:deep} gives the deep triples.
\end{proof}

\begin{corollary}[Monotone palindromic systems]\label{cor:palin}
Let $b\equiv0$, $a_n=1-c_n$, and let $0=c_0<c_1\leq c_2\leq\cdots\leq c_N=1$ satisfy $c_n+c_{N-n}=1$. Then the system has nonnegative linearization. This includes the symmetric Krawtchouk array $c_n=n/N$.
\end{corollary}
\begin{proof}
The array is palindromic. With $\sigma\equiv1$, $\alpha_n=c_n$ is nondecreasing, $\beta\equiv0$ and $\alpha_n+\gamma_n=1$. If $j\leq i$ and $i+j\leq N-2$, then $j<N-i$, so $\alpha_j=c_j\leq c_{N-i}=1-c_i=\gamma_i$. Theorem~\ref{thm:palindromic} applies.
\end{proof}

Within the palindromic subclass $b\equiv0$, $c_n+c_{N-n}=1$, the hypotheses of Corollary~\ref{cor:palin} are monotonicity inequalities. They contain the symmetric Krawtchouk endpoint, where the perturbation criterion of Section~\ref{sec:perturbation} gives no neighbourhood and where the convex counterexamples of Section~\ref{sec:open} show that convexity alone does not suffice.

Combined with the quadratic decomposition~\cite[Theorem~\ref{P1-thm:decomp}]{P1}, Corollary~\ref{cor:palin} gives an inequality criterion for symmetric systems: an array with $b\equiv0$ on $0,\ldots,2N$ has nonnegative linearization whenever its two normalized components satisfy the hypotheses of Corollary~\ref{cor:palin}. For Szwarc's system \eqref{eq:qarray} at $p_0=1/2$ the components are the symmetric Krawtchouk arrays on $N+1$ and $N$ points \cite[Corollary~\ref{P1-cor:parity}]{P1}, so this case is covered by monotonicity inequalities within the palindromic subclass.

The reversal route cannot reach systems whose diagonal increases at the top once $N\geq3$. If $b_{N-1}<b_N$, the reversed array has $g'(1,1;1)=(b'_1-b'_0)/a'_0=(b_{N-1}-b_N)/c_N<0$, and condition (b) of Theorem~\ref{thm:deep} requires $g'(1,1;1)\geq0$ when $1+1<N$. In particular it fails for the Krawtchouk arrays with $r<1$ and $N\geq3$. Szwarc's array \eqref{eq:qarray} has $b\equiv0$, so this argument does not apply to it. Its reversal replaces $p_0$ by $1-p_0$ and gives $g'(2,2;2)=(2p_0-1)/(Np_0)<0$ for $p_0<1/2$, which violates condition (b) for the array of length $2N+1$ when $4<2N$, that is, $N\geq3$. Miclo's palindromic criterion~\cite[Theorem~31]{miclo2015} requires $2^nP(n,n+1)$ to be nonincreasing on $[0,N/2]$, which the symmetric Krawtchouk chain violates for $N>2$.

\section{Limits of inequality criteria}\label{sec:limits}

Theorems~\ref{thm:low} and~\ref{thm:palindromic} control every triple except the deep ones by inequalities, the criterion of~\cite[Theorem~\ref{P1-thm:lc}]{P1} controls all triples of constant-ratio arrays with log-concave Newton profiles, and Section~\ref{sec:open} records shape conditions that fail on deep triples. This section shows that part of the failure is structural near the symmetric endpoint. For uniformly sampled polynomial deformation directions, and for the analytic and boundary-layer extensions specified below, admission at every sufficiently large symmetric Krawtchouk system imposes two first-order boundary equalities. This does not exclude arbitrary uniform inequality criteria, nor families of admissible arrays that are not of these forms. Write $K_{N,r}$ for the Krawtchouk array \eqref{eq:parameters} with $q=1$, and use the coordinates $b_1,\ldots,b_{N-1},c_1,\ldots,c_N$ of a stochastic array with $b_0=0$, $a_0=1$, $a_N=0$, $b_N=1-c_N$ and $a_n=1-b_n-c_n$ otherwise. An affine change of $x$, together with the positive normalization $R_n(x_0)=1$ at the largest support point $x_0$, reduces every finite system to this form. A direction $v=(\delta b_1,\ldots,\delta b_{N-1},\delta c_1,\ldots,\delta c_N)$ changes these coordinates; we put $\delta b_0=\delta c_0=0$ and $\delta b_N=-\delta c_N$.

\begin{lemma}[Corner value]\label{lem:corner-value}
For $0<r\leq1$, the array $K_{N,r}$ has $g(N,N;N)=(1-r)^N$.
\end{lemma}
\begin{proof}
In the model of Lemma~\ref{lem:margin}, $U_N=y_1\cdots y_N$ is the only normalized basis element containing $y_{\{1,\ldots,N\}}$, and $U_N^2=\prod_i\bigl((1-r)y_i+r\bigr)$ has coefficient $(1-r)^N$ at that monomial.
\end{proof}

Thus the corner coefficient vanishes at $r=1$, where it is a structural zero, together with the other perimeter and parity zeros of the palindromic array. The next lemma computes its first-order change exactly.

\begin{lemma}[Corner identity]\label{lem:corner}
Let $N\geq2$ and $s_N=(-1)^{\lfloor N/2\rfloor+1}\binom{N-1}{\lfloor(N-1)/2\rfloor}^{-1}$. Put $F=\delta b+2\delta c$ if $N$ is even and $F=\delta b$ if $N$ is odd. Then
\[
\frac{d}{d\varepsilon}\,g(N,N;N)\Big|_{K_{N,1}+\varepsilon v,\ \varepsilon=0}=s_N\sum_{n=0}^N(-1)^n\binom NnF(n).
\]
\end{lemma}
\begin{proof}
At $K_{N,1}$ the nodes are $\lambda_j=1-2j/N$, with masses $w_j=2^{-N}\binom Nj$ and $\pi_n=\binom Nn$ in Lemma~\ref{lem:finite}. Moreover $\binom Nn R_n(\lambda_j)$ is the binary Krawtchouk polynomial, so that
\begin{equation}\label{eq:kraw-gf}
\sum_{j=0}^N w_jR_n(\lambda_j)s^j=2^{-N}(1+s)^{N-n}(1-s)^n.
\end{equation}
The vectors $u_j=(R_n(\lambda_j))_n$ and $v_j=(w_j\pi_nR_n(\lambda_j))_n$ are right and left eigenvectors of $J$ with $v_i^\mathsf Tu_j=\boldsymbol1_{i=j}$. Since $R_N(\lambda_j)=(-1)^j$, the matrix $\Pi=R_N(J)$ reverses the indices. Let $J_\varepsilon=J+\varepsilon E$ be the perturbed recurrence matrix and $R_N^\varepsilon$ the perturbed polynomial. The matrix $R_N^\varepsilon(J_\varepsilon)$ commutes with $J_\varepsilon$ and has first row $e_N^\mathsf T$, by \eqref{eq:matrix-product}. Its derivative $X$ at $\varepsilon=0$ therefore satisfies $[J,X]=[\Pi,E]$ and $e_0^\mathsf TX=0$. In the eigenbasis, $\widehat X_{ij}=v_i^\mathsf TXu_j$ equals $((-1)^i-(-1)^j)\widehat E_{ij}/(\lambda_i-\lambda_j)$ for $i\neq j$, where $\widehat E_{ij}=v_i^\mathsf TEu_j$, and $\sum_i\widehat X_{ij}=0$. By \eqref{eq:matrix-product}, and since $\pi_N=1$, the derivative in the lemma is $h_N:=e_N^\mathsf TXe_N=\sum_{i,j}(-1)^{i+j}\widehat X_{ij}w_j$. Eliminating $\widehat X_{jj}$ gives
\[
h_N=-2\sum_{i+j\ \mathrm{odd}}\widehat X_{ij}w_j=N\sum_{i\neq j}\widehat E_{ij}w_j\frac{(-1)^i-(-1)^j}{i-j}.
\]
Perturbing row $n$ alone means $E=e_n\rho^\mathsf T$ and $\widehat E_{ij}=w_i\pi_nR_n(\lambda_i)\,\rho^\mathsf Tu_j$. Write $\mathrm i$ for the imaginary unit and insert
\[
\frac{(-1)^i-(-1)^j}{i-j}=-\pi\int_0^1\sin\pi\bigl(\theta i+(1-\theta)j\bigr)\,d\theta,
\]
which also vanishes for $i=j$. Then
\begin{gather*}
h_N=-\pi N\pi_n\int_0^1\operatorname{Im}\bigl(A(\theta)B(\theta)\bigr)\,d\theta,\\
A=\sum_iw_iR_n(\lambda_i)e^{\mathrm i\pi\theta i},\qquad
B=\sum_jw_j(\rho^\mathsf Tu_j)e^{\mathrm i\pi(1-\theta)j},
\end{gather*}
and both sums are given by \eqref{eq:kraw-gf}. Put $\phi=\pi\theta/2$.
\begin{itemize}
\item Coordinate $b_n$, $1\leq n<N$, so $\rho=e_n-e_{n+1}$. Then
\[
AB=-(-1)^n\,\mathrm i^{N-1}\cos^N\!\phi\,\sin^{N-1}\!\phi\,e^{-\mathrm i\phi},
\]
and $h_N=s_N(-1)^n\binom Nn$. The integral reduces to $\int_0^{\pi/2}\cos^{N+1}\!\phi\,\sin^{N-1}\!\phi\,d\phi$ ($N$ even) or $\int_0^{\pi/2}\cos^{N}\!\phi\,\sin^{N}\!\phi\,d\phi$ ($N$ odd), and each equals $|s_N|/(2N)$.
\item Coordinate $c_n$, $1\leq n<N$, so $\rho=e_{n-1}-e_{n+1}$. Now $B$ acquires the factor $e^{\mathrm i\phi}/\cos\phi$, and the integrand carries $\sin((N-1)\pi/2)$. The contribution vanishes for odd $N$. For even $N$ it is twice that of $b_n$, since the two beta integrals have ratio $2$.
\item Coordinate $c_N$, so $\rho=e_{N-1}-e_N$ in row $N$. The same computation with $n=N$ gives $s_N$.
\end{itemize}
Summing these contributions with the conventions $\delta b_0=\delta c_0=0$ and $\delta b_N=-\delta c_N$ gives the stated alternating sum.
\end{proof}

\begin{definition}\label{def:rule}
A \emph{uniform rule} is a pair of real functions $\beta,\gamma$ on $[0,1]$. It defines, for each $N$, the direction $v_N$ with $\delta b_n=\beta(n/N)$ for $1\leq n<N$ and $\delta c_n=\gamma(n/N)$ for $1\leq n\leq N$. A set of arrays \emph{admits} the rule at the symmetric endpoint if there is $N_0$ such that for every $N\geq N_0$ there is $\varepsilon_N>0$ with $K_{N,1}+\varepsilon v_N$ in the set for all $0<\varepsilon<\varepsilon_N$. Put
\[
L_{\mathrm{even}}=\beta(0)+2\gamma(0)+\beta(1)+\gamma(1),\qquad L_{\mathrm{odd}}=\beta(0)-\beta(1)-\gamma(1).
\]
\end{definition}

\begin{theorem}[Forced equalities for uniform criteria]\label{thm:uniform}
Let $\beta,\gamma$ be polynomials. If a set of arrays with property (P) admits the rule $(\beta,\gamma)$ at the symmetric endpoint, then $L_{\mathrm{even}}=L_{\mathrm{odd}}=0$. If $L_{\mathrm{even}}=L_{\mathrm{odd}}=0$, the first-order change of $g(N,N;N)$ along $v_N$ vanishes for every $N$ above the degrees of $\beta$ and $\gamma$.
\end{theorem}
\begin{proof}
For a polynomial $p$ of degree less than $N$, $\sum_{n=0}^N(-1)^n\binom Nnp(n/N)=0$. Only boundary terms survive in Lemma~\ref{lem:corner}: the derivative equals $-s_NL_{\mathrm{even}}$ for even $N$ and $-s_NL_{\mathrm{odd}}$ for odd $N$. The sign of $s_N$ changes from $N$ to $N+2$. If, for instance, $L_{\mathrm{even}}\neq0$, then the derivative is negative for every large even $N$ in one residue class modulo $4$. For such $N$, $g(N,N;N)<0$ at $K_{N,1}+\varepsilon v_N$ for all small $\varepsilon>0$, since the coefficient vanishes at $\varepsilon=0$.
\end{proof}

The known positive deformations satisfy both equalities: the Krawtchouk direction $\beta(x)=x$, $\gamma(x)=-x$, palindromic deformations with $\beta=0$, $\gamma$ odd about $\tfrac12$ and $\gamma(0)=\gamma(1)=0$ (Corollary~\ref{cor:palin}), and the tangent direction $\beta=0$, $\gamma(x)=x(x-1)$ of the geometric family at $q=1$ (Theorem~\ref{thm:criterion}). The two functionals are independent already on constant rules, where $L_{\mathrm{even}}=2\beta+3\gamma$ and $L_{\mathrm{odd}}=-\gamma$. Consequently, for every $d\geq0$, no set of arrays with property (P) admits all rules of a nonempty open subset of the polynomial rules of degree at most $d$. In this sense, a criterion that admits the symmetric Krawtchouk systems through polynomial uniform rules must impose the two equalities on them. The geometric criterion and Corollary~\ref{cor:palin} are of this mixed type.

The polynomial hypothesis can be replaced by analyticity on a fixed region. Let $U(x)=\int_0^1\log|x-t|\,dt$ and $\Omega=\{x\in\mathbb C:U(x)\leq-1\}$, the region of convergence of equispaced interpolation. It contains $[0,1]$, has nonanalytic boundary points (with vertical tangents) at $0$ and $1$, and meets the line $\operatorname{Re}x=\tfrac12$ at $\operatorname{Im}x=\pm0.26276\ldots$.

\begin{lemma}[N\"orlund--Rice bound]\label{lem:nr}
If $f$ is analytic on an open set containing $\Omega$, then $\sum_{n=0}^N(-1)^n\binom Nnf(n/N)=O(\sqrt N\,e^{-\eta N})$ for some $\eta>0$.
\end{lemma}
\begin{proof}
The residues of $1/\prod_{k=0}^N(x-k/N)$ give
\[
\sum_{n=0}^N(-1)^n\binom Nnf(n/N)=(-1)^N\frac{N!}{N^N}\,\frac1{2\pi i}\oint_\Gamma\frac{f(x)\,dx}{\prod_{k=0}^N(x-k/N)}
\]
for any cycle $\Gamma$ in the domain of $f$ that winds once around $[0,1]$. Since $U$ is continuous and tends to $\infty$, the compact sets $\{U\leq-1+t\}$ shrink to $\Omega$ as $t\downarrow0$. Choose $\eta>0$ with $\{U\leq-1+3\eta\}$ inside the domain of $f$, and a bounded polygonal open set $V$ with $\{U\leq-1+\eta\}\subset V$ and $\overline V\subset\{U<-1+3\eta\}$. The oriented boundary $\Gamma$ of $V$, possibly with several components, is a cycle outside $\{U\leq-1+\eta\}$ that winds once around every node $k/N$ and bounds a region inside the domain of $f$. With $d=\operatorname{dist}(\Gamma,[0,1])>0$, the derivative of $t\mapsto\log|x-t|$ on $[0,1]$ is bounded by $1/d$ for $x\in\Gamma$, so a Riemann-sum estimate gives $\sum_{k=0}^N\log|x-k/N|\geq NU(x)-C_\Gamma\geq N(-1+\eta)-C_\Gamma$ with $C_\Gamma$ independent of $N$. Stirling's bound $N!/N^N\leq e\sqrt Ne^{-N}$ completes the proof.
\end{proof}

\begin{corollary}\label{cor:analytic-rules}
Let $\beta,\gamma$ be analytic on an open set containing $\Omega$. If a set of arrays with property (P) admits the rule $(\beta,\gamma)$ at the symmetric endpoint, then $L_{\mathrm{even}}=L_{\mathrm{odd}}=0$. When these equalities hold, the first-order change of $g(N,N;N)$ along $v_N$ is $O(|s_N|\sqrt N\,e^{-\eta N})$; it need not vanish.
\end{corollary}
\begin{proof}
By Lemma~\ref{lem:nr}, the alternating sum of $\beta+2\gamma$ or $\beta$ is $O(\sqrt Ne^{-\eta N})$. The derivative in Lemma~\ref{lem:corner} is therefore $s_N(-L_{\mathrm{even}}+o(1))$ for even $N$ and $s_N(-L_{\mathrm{odd}}+o(1))$ for odd $N$, and the argument of Theorem~\ref{thm:uniform} gives the equalities. The rule $\beta(x)=e^x-1-(e-1)x$, $\gamma=0$ has $L_{\mathrm{even}}=L_{\mathrm{odd}}=0$ and derivative $s_N(1-e^{1/N})^N\neq0$.
\end{proof}

Deformations may also depend on $N$ in other ways. For each fixed $N$ the symmetric endpoint is a limit of systems whose product coefficients are strictly positive on the whole triangle; the corner nevertheless constrains families that are regular in $N$. The first statement rests on the following lemma, valid for every finite system. Write $\mu$ for the orthogonality measure and $x_0$ for its largest support point, so that $R_n(x_0)=1$.

\begin{lemma}[Lowering the top mass]\label{lem:topmass}
Let $0\leq\varepsilon<\mu(\{x_0\})$, let $\mu_\varepsilon=\mu-\varepsilon\delta_{x_0}$, and let $R_n^\varepsilon$ be its polynomials, normalized at $x_0$. Put $\pi_n^\varepsilon=\bigl(\int(R_n^\varepsilon)^2\,d\mu_\varepsilon\bigr)^{-1}$ and $\pi_n=\pi_n^0$; these inverse norms are not rescaled to impose $\pi_0^\varepsilon=1$. The recurrence coefficients in~(1) are taken in the fixed coordinate $x$, without an $\varepsilon$-dependent affine normalization.
\begin{enumerate}
\item The recurrence coefficients change at rate $(\dot a_n,\dot b_n,\dot c_n)=\pi_n(a_n,\,c_n-a_n,\,-c_n)$.
\item If $\int R_xR_yR_z\,d\mu=0$, then
\begin{gather*}
\frac{d}{d\varepsilon}\Big|_{\varepsilon=0}\int R^\varepsilon_xR^\varepsilon_yR^\varepsilon_z\,d\mu_\varepsilon=2-T_x(y,z)-T_y(x,z)-T_z(x,y),\\
T_x(y,z)=\sum_{k\geq x}g(y,z;k).
\end{gather*}
\end{enumerate}
\end{lemma}
\begin{proof}
Write $R_n^\varepsilon=R_n+\varepsilon Q_n+O(\varepsilon^2)$, with $\deg Q_n\leq n$ and $Q_n(x_0)=0$. First-order orthogonality gives $\int Q_np\,d\mu=p(x_0)$ for $\deg p<n$, hence $Q_n=\sum_{k<n}\pi_k(R_k-R_n)$, and $\pi_n^\varepsilon=\pi_n\bigl(1+\varepsilon(2S_n+\pi_n)\bigr)+O(\varepsilon^2)$ with $S_n=\sum_{k<n}\pi_k$. For (1), differentiate $c_n=\pi_{n-1}\int xR_nR_{n-1}\,d\mu$ and $b_n=\pi_n\int xR_n^2\,d\mu$, using the recurrence, the row sums $a_m+b_m+c_m=x_0$ and $\pi_{n-1}a_{n-1}=\pi_nc_n$; the terms with $S_n$ cancel, and $\dot a_n$ follows from the row sums. For (2), the derivative equals $-R_xR_yR_z(x_0)+\int(Q_xR_yR_z+R_xQ_yR_z+R_xR_yQ_z)\,d\mu$, and $\int Q_xR_yR_z\,d\mu=\sum_{k<x}g(y,z;k)=1-T_x(y,z)$.
\end{proof}

\begin{proposition}[Interior points near the symmetric endpoint]\label{prop:interior}
For every $N\geq2$ there is $\Lambda_N<\infty$ such that, for $\Lambda>\Lambda_N$ and all sufficiently small $\varepsilon>0$, the measure of $K_{N,1-\Lambda\varepsilon}$ with its top mass lowered by $\varepsilon$ has $g(m,n;k)>0$ for all $|m-n|\leq k\leq\min(m+n,N)$. Consequently, after the affine normalization of Section~\ref{sec:criterion}, a relatively open set of arrays near it has property (P).
\end{proposition}
\begin{proof}
By the model of Lemma~\ref{lem:margin}, $K_{N,r}$ has
\begin{gather*}
g(m,n;k)=\sum_jH(m,n;j)\binom js r^{j-s}(1-r)^s,\\
H(m,n;j)=\frac{\binom mj\binom{N-m}{n-j}}{\binom Nn},\qquad s=k-m-n+2j,
\end{gather*}
where the sum is over $\max(0,m+n-N)\leq j\leq\min(m,n)$ with $0\leq s\leq j$.
At $r=1$ only $s=0$ survives, so a triangle coefficient of $K_{N,1}$ vanishes exactly when $m+n+k$ is odd or $m+n+k>2N$. Both deformations are smooth in $(r,\varepsilon)$, so only these zeros need first-order positivity.
\begin{itemize}
\item If $m+n+k>2N$, all three tails in Lemma~\ref{lem:topmass}(2) vanish, since $g(y,z;k)=0$ for $k>2N-y-z$. The top-mass rate is therefore $2\pi_k>0$, and the rate of decreasing $r$ is nonnegative because $g\geq0$ for $r\leq1$.
\item If $m+n+k$ is odd and at most $2N+1$, the rate of decreasing $r$ is the $s=1$ term $H(m,n;j)\,j$ with $j=(1+m+n-k)/2$. It is positive, because $j\geq1$ and $\max(0,m+n-N)\leq j\leq\min(m,n)$.
\end{itemize}
These cases cover all zeros. Let $\mathcal Z_N$ be the set of triangle triples whose coefficient vanishes at $K_{N,1}$. For $z\in\mathcal Z_N$, write $A_z$ for the top-mass rate and $B_z$ for the rate of decreasing $r$, and define
\[
\Lambda_N=\max\Bigl(\{0\}\cup\Bigl\{-\frac{A_z}{B_z}:z\in\mathcal Z_N,\ A_z<0\Bigr\}\Bigr).
\]
Every denominator appearing here is positive by the preceding cases. Hence $\Lambda>\Lambda_N$ makes every reference-zero derivative strictly positive, and finiteness of $\mathcal Z_N$ gives a common threshold for $\varepsilon$.

Here the measure of $K_{N,r}$ is its probability orthogonality measure. Choose $\varepsilon$ so small that $r=1-\Lambda\varepsilon>0$ and $\varepsilon<(r/(1+r))^N$, so that the lowered top mass remains positive. The affine coordinate $y=(1-\varepsilon)x+\varepsilon=R_1^\varepsilon(x)$ restores the basic recurrence normalization and leaves all product coefficients unchanged. Strict positivity on the triangle gives $b_n=g(1,n;n)>0$ for $1\leq n\leq N$. Lemma~\ref{lem:stability}, with $\Sigma$ the whole triangle, therefore gives a relatively open neighborhood with property (P) in the affine space defined by the basic boundary and row-sum equalities.
\end{proof}

Exact rational computation of the maximum above gives $\Lambda_N=4,10,18,28,52,90,165$ for $N=4,\ldots,10$. These finite-size values do not establish an asymptotic growth law or a geometric thickness estimate. At $K_{N,1}$, in the fixed spectral coordinate, Lemma~\ref{lem:topmass}(1) gives
\[
\dot c_n=-\frac nN\binom Nn,\qquad\dot b_n=\Bigl(\frac{2n}{N}-1\Bigr)\binom Nn.
\]
After the affine normalization restoring $R_1=x$, the top-mass direction is
\[
\delta a_n=\frac{N-n}{N}\Bigl(\binom Nn-1\Bigr),\qquad
\delta b_n=1+\Bigl(\frac{2n}{N}-1\Bigr)\binom Nn,\qquad
\delta c_n=-\frac nN\Bigl(1+\binom Nn\Bigr).
\]
Its binomial dependence on $N$ is not a uniform rule of Definition~\ref{def:rule}, consistently with Theorem~\ref{thm:uniform}.

\begin{proposition}[Families depending on $N$]\label{prop:families}
Let $v_N$ be directions such that a set with property (P) contains $K_{N,1}+\varepsilon v_N$ for all large $N$ and small $\varepsilon>0$, and let $A_N$ be the alternating sum in Lemma~\ref{lem:corner} for $v_N$. Then $s_NA_N\geq0$ for all large $N$.
\begin{enumerate}
\item Along either parity class, $A_N$ is not eventually strictly positive and not eventually strictly negative. If $A_N$ has an eventual weak sign along a parity class, it vanishes for all large $N$ of the residue class modulo $4$ on which that sign conflicts with the sign of $s_N$.
\item Suppose $\delta b_n=\beta(n/N)+p_n+q_{N-n}$ and $\delta c_n=\gamma(n/N)+p'_n+q'_{N-n}$, with $\beta,\gamma$ as in Corollary~\ref{cor:analytic-rules} and finitely supported boundary layers $p,q,p',q'$ independent of $N$. Then $q'_0=L_{\mathrm{even}}=L_{\mathrm{odd}}$, $p_k=q_k$ and $p'_k+q'_k=-p_k$ for $k\geq1$.
\end{enumerate}
\end{proposition}
\begin{proof}
A negative first-order change of $g(N,N;N)$ excludes $K_{N,1}+\varepsilon v_N$ for small $\varepsilon$, so $s_NA_N\geq0$. Part (1) follows because the sign of $s_N$ alternates along each parity class. In (2), write $D_Nf=\sum_{n=0}^N(-1)^n\binom Nnf(n/N)$. For large $N$, $A_N=D_N(\beta+2\gamma)+P_e(N)$ for even $N$ and $A_N=D_N\beta+P_o(N)$ for odd $N$, where
\begin{gather*}
P_e(N)=q'_0-L_{\mathrm{even}}+\sum_{k\geq1}(-1)^k\binom Nk\bigl[p_k+q_k+2(p'_k+q'_k)\bigr],\\
P_o(N)=q'_0-L_{\mathrm{odd}}+\sum_{k\geq1}(-1)^k\binom Nk(p_k-q_k).
\end{gather*}
These are polynomials in $N$. A nonzero polynomial has an eventual strict sign, which the remainder $O(\sqrt Ne^{-\eta N})$ of Lemma~\ref{lem:nr} cannot change, so part (1) forces $P_e\equiv P_o\equiv0$. The polynomials $1,\binom N1,\binom N2,\ldots$ are linearly independent, which gives the stated conditions.
\end{proof}

Thus a separate adjustment of the top coefficient $c_N$ reduces the two bulk equalities to the single equality $\gamma(0)+\beta(1)+\gamma(1)=0$. For $k\geq1$ and $N\geq k+K+1$, where $K$ bounds the supports of $p,q,p',q'$, the boundary-layer contributions satisfy the palindromic first-order relations $\delta b_k=\delta b_{N-k}$ and $\delta a_k=\delta c_{N-k}$; the separate terminal adjustment need not itself be palindromic. Families not covered by part~(1) must have $A_N$ change sign, or vanish, infinitely often within a parity class.

Criteria that cover $r<1$ but not $r=1$ are thin near $r=1$.

\begin{proposition}[Thinness near the symmetric endpoint]\label{prop:thin}
Fix $N\geq2$ and a direction $v$ along which the first-order change $\dot g$ of $g(N,N;N)$ at $K_{N,1}$ is negative. There are $r_*<1$ and $C,c>0$ such that $K_{N,r}+\varepsilon v$ lacks property (P) whenever $r_*<r<1$ and $C(1-r)^N\leq\varepsilon\leq c$.
\end{proposition}
\begin{proof}
The function $G(r,\varepsilon)=g(N,N;N)$ at $K_{N,r}+\varepsilon v$ is a rational function of the recurrence entries whose denominators involve only upward coefficients, which are positive at $K_{N,1}$; so $G$ is smooth near $(1,0)$. By Lemma~\ref{lem:corner-value}, $G(r,0)=(1-r)^N$. Put $d=|\dot g|$. On a small rectangle around $(1,0)$, on which the basic sign conditions persist, $\partial_\varepsilon G\leq-d/2$, hence $G(r,\varepsilon)\leq(1-r)^N-d\varepsilon/2<0$ for $\varepsilon\geq4(1-r)^N/d$.
\end{proof}

\begin{corollary}[Optimal order of the stability radius]\label{cor:radius-order}
For $N\geq2$ let $\rho_N(r)$ be the largest $\rho$ such that every array satisfying \eqref{eq:basic} with $\widetilde b_i\geq0$ and with left side of \eqref{eq:closeness} at most $\rho$ has property (P). Then $\rho_N(r)$ is of exact order $(1-r)^N$ as $r\to1$.
\end{corollary}
\begin{proof}
Theorem~\ref{thm:perturbation} gives $\rho_N(r)\geq\eta_N(r)=\tfrac12((1-r)/(32N))^N$ for $r\geq\tfrac12$. For the upper bound, use the terminal direction $\delta c_N=-\operatorname{sgn}(s_N)$, $\delta b_N=\operatorname{sgn}(s_N)$, with all other entries unchanged. By Lemma~\ref{lem:corner} its first-order change is $-|s_N|<0$, and it moves the full array by exactly $\varepsilon$ in the norm of \eqref{eq:closeness}. Put $t=1-r$ and $\varepsilon=Ct^N$ as in Proposition~\ref{prop:thin}. For $N\geq2$ and $r$ near $1$, $\varepsilon\leq t/2$ and $\varepsilon\leq r/2$, so $\widetilde b_N\geq t/2$ and $\widetilde c_N\geq r/2$; every other coefficient is unchanged. The array satisfies the basic conditions and $\widetilde b_i\geq0$, lies at full-array distance $\varepsilon=C(1-r)^N$ from $K_{N,r}$, and lacks property (P). Hence $\rho_N(r)\leq C(1-r)^N$. The restriction $N\geq2$ is needed: every admissible two-point array with $b_1\geq0$ has property (P).
\end{proof}


For fixed $r$ the radius also decays exponentially in $N$: raising the top mass of the measure of $K_{N,r}$, the reverse of Lemma~\ref{lem:topmass}, gives $\rho_N(r)\leq3(1-r^2)^N$ whenever $(1-r^2)^N\leq\frac12$ and $3(1-r^2)^N\leq(1-r)/N$~\cite[Proposition~\ref{P1-prop:raise} and Corollary~\ref{P1-cor:radius-exp}]{P1}.

Thus, for every fixed $r\in(0,1)$, no set of arrays with property (P) contains coefficient balls of a positive radius independent of $N$ around all $K_{N,r}$; Theorem~\ref{thm:perturbation} gives the lower bound $\rho_N(r)\geq\tfrac12(\min(r,1-r)/(32N))^N$. The deformation that raises the top mass has the binomial profile of the masses and is not a uniform rule. Along the terminal row alone, the first-order change of $g(N,N;N)$ at $K_{N,r}$ has an exact contour-integral representation, obtained as in Lemma~\ref{lem:corner} with $R_N(J)$ having eigenvalues $(-r)^j$. The formal saddle contributions have exponential factor $(1-r+r^2)/(1+r)$, which exceeds $1-r$ exactly when $r>\tfrac12$. For $r>\tfrac12$ this suggests an oscillatory first-order envelope, not a proved threshold for finite perturbations. No asymptotic estimate or uniform control of the nonlinear remainder is proved here. Exact computations (\path{fixedr_probe.py}) confirm failures of property (P) along the terminal row at sizes down to $10^{-15}$ for $r=4/5$, $N=40$. The comparison at $r=\tfrac12$ concerns only the first-order term. Along the terminal row, $g(N,N;N)$ is a polynomial in the perturbation, since the recurrence for the linearization coefficients divides only by $a_0,\ldots,a_{N-1}$.  Exact terminal-row failures also occur for $r\leq\tfrac12$: at size $6.3\cdot10^{-6}$ for $r=\tfrac12$, $N=120$; at size $3.2\cdot10^{-3}$ for $r=\tfrac25$, $N=100$; and at size $1.6\cdot10^{-2}$ for $r=\tfrac13$, $N=120$.

\paragraph{Computations away from $r=1$.}
The explicit fixed-size, one-kink, and log-concave-increment examples below are reproduced by \path{rigidity_probe.py} and \path{smallr_probe.py}. Exact rational computations certify the stated finite examples; the separately labelled high-precision calculations provide numerical evidence. No asymptotic conclusion is asserted. Here $g(N,N;N)$ is computed from column $N$ of the product table, an $O(N^2)$-operation recurrence, so sizes in the hundreds are cheap.
\begin{itemize}
\item For the rule $\beta=1$, $\gamma=0$, exact computation shows that property (P) fails at $r=4/5$, with perturbation size $1/100$, for every $N=6,\ldots,14$, and at $r=19/20$, with perturbation size $1/1000$, for every $N=4,\ldots,14$. The period-four sign pattern at $r=1$ does not persist in general for $r<1$: at $r=4/5$ the first-order derivatives are $512/1875$ for $N=4$ and $-71936/5859375$ for $N=8$.
\item Below $r=\tfrac12$, the counterexamples we found to shape conditions occur at sizes that grow as $r$ decreases; we did not determine minimal sizes or prove a scaling. For constant-ratio arrays with convex profile, one-kink counterexamples ($m$ unit increments, then $1+\delta$) have $g(N,N;N)<0$ in exact arithmetic at
\begin{center}
\begin{tabular}{@{}lccccccc@{}}
$r$&$2/5$&$1/3$&$3/10$&$1/4$&$2/7$&$1/5$&$1/10$\\
$N$&$20$&$40$&$40$&$60$&$80$&$100$&$500$
\end{tabular}
\end{center}
\item Convex profiles with log-concave increments (step ratios nonincreasing and at least $1$) contain Krawtchouk, every geometric profile and the normalized quadratic profiles $d_\kappa(x)=(x+\kappa x^2)/(1+\kappa)$, $\kappa\geq0$, and they exclude the one-kink counterexamples. Yet $k$ steps of ratio $1+1/(2N)$ followed by flat increments give $g(N,N;N)<0$ at $(r,N,k)=(\tfrac12,80,40)$, $(\tfrac25,32,22)$ and $(\tfrac13,90,62)$ in exact arithmetic. In $80$- and $150$-digit floating-point arithmetic, which agree but are not exact sign certificates, the same happens at $(\tfrac14,250,175)$ and $(\tfrac15,500,400)$.
\end{itemize}
The exact examples above disprove the stated shape conditions at the listed parameters and sizes. The two larger high-precision examples provide additional numerical evidence. These findings establish neither minimal failure sizes nor general conclusions about smooth deformations.

\section{Open problems}\label{sec:open}

\paragraph{Sharper stability radii.}
For fixed $N$ and $0<r<1$, determine or estimate the largest coefficient radius for which every admissible stochastic perturbation has property (P). Corollary~\ref{cor:radius-order} shows that, for fixed $N\geq2$, the order $(1-r)^N$ as $r\to1$ is optimal, and~\cite[Corollary~\ref{P1-cor:radius-exp}]{P1} gives $\rho_N(r)=O((1-r^2)^N)$ for each fixed $r$; the gap to the lower bound $\tfrac12(\min(r,1-r)/(32N))^N$ of Theorem~\ref{thm:perturbation} remains open. The bound in Theorem~\ref{thm:perturbation} comes from a single-term lower margin and a worst-case recurrence amplification. Both estimates discard substantial information. A useful improvement would retain a transparent coefficient test and describe its dependence on both $N$ and the bias parameter.

\paragraph{The symmetric endpoint.}
At $r=1$, the reference margin in Lemma~\ref{lem:margin} vanishes because parity forces many coefficients to be zero. Symmetric perturbations preserve the parity zeros, and Theorem~\ref{thm:symmetric} exploits this for Szwarc's quadratic system when $p_0<1/2$. At $p_0=1/2$, and for the symmetric Krawtchouk array, further coefficients on the parity support vanish. Determine perturbation directions and structural constraints that preserve positivity there; continuity alone does not control coefficients emerging from those zeros. For every $N\geq1$, the symmetric Krawtchouk array is a limit of arrays whose product coefficients are strictly positive throughout the triangle support. For $N\geq2$ this follows from Proposition~\ref{prop:interior}; the two-point case is immediate. Theorem~\ref{thm:uniform} and Corollary~\ref{cor:analytic-rules} impose two boundary equalities on admitted polynomial rules and on the analytic rules satisfying the holomorphic-extension hypothesis specified there. A description of the cone for all $N$ is open.

\paragraph{Monotonicity-type criteria for finite systems.}
Szwarc's Theorems~3.2 and~3.3 in~\cite{szwarc} are open conditions of monotonicity type for infinite systems. For finite systems, property (P) at the boundary point is the hypergroup property of the corresponding birth--death chain, and Miclo~\cite[Theorem~31 and Proposition~6]{miclo2015} adapts the wave-equation argument of Bakry and Huet to give inequality conditions for chains that are symmetric under $n\mapsto N-n$ and sufficiently convex. These conditions do not cover the full Krawtchouk or Szwarc families. For symmetric Krawtchouk they fail when $N>2$, but the cases $N=1,2$ are covered; the Szwarc instance $N=1$, $p_0=1/2$ is the same covered three-point system. By the quadratic decomposition~\cite[Theorem~\ref{P1-thm:decomp}]{P1}, any criterion of this kind for nonsymmetric finite systems, valid including the terminal indices, immediately gives one for symmetric finite systems. Theorem~\ref{thm:low} supplies such hypotheses for all triples except those whose pairwise sums all exceed $N$. Theorem~\ref{thm:all} supplies such hypotheses for all triples under its strengthened conditions (i)--(iii), and Corollary~\ref{cor:heavy} contains the arrays \eqref{eq:parameters} with $q\geq2$ but not the Krawtchouk arrays with $N\geq3$. For constant-ratio arrays, the log-concave Newton criterion~\cite[Theorem~\ref{P1-thm:lc}]{P1} controls the deep triples by inequalities that contain the whole geometric criterion, in particular the Krawtchouk arrays, and~\cite{P2} extends it to ratios $c_i/(b_i+c_i)$ that increase in several ways; other nonconstant ratios remain open. Convexity of the profile does not suffice even with a constant ratio and $r$ bounded away from $1$: for $d$ with increments $1$ at the first seven steps and $6/5$ at the remaining seven ($N=14$), exact arithmetic gives $g(14,14;14)<0$ at $r=1/2$, and the analogous profile with $N=20$, twelve unit increments and $21/20$ afterwards gives $g(20,20;20)<0$ at $r=2/5$. Log-convex increments fail as well: fourteen unit increments followed by $(21/20)^t$, $t=1,\ldots,4$, give $g(18,18;18)<0$ at $r=1/2$. These examples show that the stated shape restrictions do not guarantee positivity at the specified biases and sizes. Each failure is witnessed by a negative terminal coefficient $g(N,N;N)$. Section~\ref{sec:limits} explains part of this: admission at the symmetric endpoint imposes two first-order boundary equalities for the polynomial rules of Theorem~\ref{thm:uniform} and the analytic rules of Corollary~\ref{cor:analytic-rules}, and the computational discussion there gives counterexamples to convex-profile conditions and to convex-profile conditions with log-concave increments at the stated parameters and sizes. Two questions remain. First,~\cite[Corollary~\ref{P1-cor:radius-exp}]{P1} gives an exponential upper bound on the admissible coefficient-ball radius around $K_{N,r}$ for every fixed $r\in(0,1)$, witnessed by raising the top mass. Determine whether analogous upper bounds hold for terminal-row perturbations, and which uniform rules exhibit such thinness. The first-order comparison in the remark after Corollary~\ref{cor:radius-order} changes at $r=\tfrac12$. The exact terminal-row failures reported there for $r\leq\tfrac12$ show that higher-order terms matter, so this comparison does not locate a threshold for finite perturbations. Second, decide whether a criterion given by linear inequalities in the coefficients can admit every $K_{N,r}$ and $G_{N,q,r}$ ($N\geq1$, $q\geq1$) for all $r$ in a range $0<r\leq r_0$. A criterion given by linear inequalities in the coefficients admits a convex set of arrays, but the set of arrays satisfying property (P) need not be convex, even at fixed $(N,r)$ and with both endpoints admitted by Theorem~\ref{thm:criterion}. For $q\geq1$ let $G_{N,q,r}$ be the array \eqref{eq:parameters}, so that $G_{N,1,r}=K_{N,r}$. For $q>1$ and $0<\lambda<1$ the combination $(1-\lambda)K_{N,r}+\lambda G_{N,q,r}$ has increments proportional to $1+t\,q^{i-N}$, where $t=\lambda N(q-1)q^{N-1}/((1-\lambda)(q^N-1))$. These combinations have $g(N,N;N)<0$ at seven exactly certified parameter choices with $r$ between $2/7$ and $4/5$~\cite[Section~\ref{P1-sec:limits}]{P1}, and, in $150$- and $300$-digit arithmetic that is not an exact sign certificate, at $(r,N,q,t)=(1/4,500,3,10)$. At the certified parameters no sufficient criterion whose admitted set is convex, in particular none defined by a system of linear inequalities in the coefficients, can admit both $K_{N,r}$ and $G_{N,q,r}$. In our searches, the sizes of the examples found tended to increase as $r$ decreased; at $r=1/5$ no negative corner was found on the tested $(N,q,t)$ grid with $N\leq1000$. Sign conditions on the shape fail in the same way. Write $d_n=b_n+c_n$ and $\rho_n=c_n/d_n$ for $n\geq1$. For $S_1,S_2\in\{[0,\infty),(-\infty,0],\{0\},\mathbb R\}$ and $\rho_{\max}>0$, consider the stochastic arrays with $b_0=0$, $a_n$ decreasing, $b_n$ and $c_n$ nondecreasing, $d_{n+2}-2d_{n+1}+d_n\in S_1$ for $0\leq n\leq N-2$, $\rho_{n+1}-\rho_n\in S_2$ for $1\leq n\leq N-1$, and $\rho_n\leq\rho_{\max}$. Each such class contains every $K_{N,r}$ with $r\leq\rho_{\max}$, and for $S_1=S_2=\{0\}$ it consists of these arrays only. For every other choice of $(S_1,S_2)$ and every $\rho_{\max}\geq83/400$, exact computation (\path{signpattern_probe.py}) gives an array in the class with $g(N,N;N)<0$. In the first two examples, normalize the listed increments to sum to one. If $S_1$ contains a positive number, take constant ratio $1/5$ and eighty unit increments followed by twenty increments $11/10$ ($N=100$). If $S_1$ contains a negative number, take constant ratio $1/10$ and seventeen unit increments followed by three increments $1/2$ ($N=20$). If $S_1=\{0\}$ and $S_2$ contains a positive number, take $d_n=n/N$ and $\rho_n=1/5+\max(n-102,0)/2400$ ($N=120$). If $S_1=\{0\}$ and $S_2$ contains a negative number, take $d_n=n/N$ and $\rho_n=1/5-\max(n-68,0)/1600$ ($N=80$). Thus, for every $\rho_{\max}\geq83/400$, among the classes just defined, only $S_1=S_2=\{0\}$ guarantees property (P) for all $N$. For $\rho_{\max}<83/400$, these four examples do not settle all sign choices; the cases not excluded by them remain unresolved here. For instance, with $S_1=\{0\}$ and increasing ratio starting from $\rho=1/10$, our search found no failure for $N\leq300$. Criteria with conditions on higher differences, with mixed conditions, or with conditions that do not apply at every index are not covered. Floating-point searches for Szwarc-type certificates on the deep triples, which are not exact, are recorded in the companion code.

\paragraph{Other finite families.}
Find analogous coefficient regions around other finite orthogonal systems and compare their geometry with the Krawtchouk neighborhoods. An algebra with a positive basis may supply the reference margin, but a recurrence criterion must express the resulting admissible deformations without requiring the unknown product coefficients as input.


\section*{AI-use disclosure}

OpenAI Codex was used to develop the proof arguments of Sections~\ref{sec:criterion}--\ref{sec:applications}, search and check source passages, construct exact-arithmetic diagnostics, draft and revise this manuscript, and prepare the Lean translation. Separate AI-assisted reviews checked the algebraic proof, boundary cases, normalizations, and quantitative estimates; these reviews are not human peer review. Anthropic Claude was used to analyse Szwarc's Problem~2 against the survey text, to explore discrete hyperbolic certificates numerically, and to develop and draft Sections~\ref{sec:szwarc}--\ref{sec:limits} together with their exact diagnostics, and it prepared the Lean formalizations of those sections. These sections were reviewed in fourteen rounds by ChatGPT Pro, as part of a longer manuscript (an AI review with independent exact checks, not human peer review); the first round failed, the corrections are recorded in \path{REVIEW.md}, and for each section the final confirmation round on the corrected text returned PASS. The author set the research objective and retains responsibility for publication decisions and the final mathematical claims. No independent human verification is asserted here. ``Machine-verified'' refers only to the Lean theorems listed in the section on code and formalization, with the axioms stated there. Results and prose remarks not covered by those listed Lean theorems are supported by written proofs, AI review and exact numerical checks, but are not machine-verified.

\section*{Code and formalization}

The companion code~\cite{code} contains exact-arithmetic diagnostics for every result (\path{papers/DIAGNOSTICS.md}), among them \path{verify.py}, which compares recurrence-matrix multiplication with multiplication modulo the terminal polynomial, and a Lean~4 development, built and audited by \verb|make lean| and \verb|make axioms|; the audit rejects dependencies outside \path{propext}, \path{Classical.choice} and \path{Quot.sound}. All declarations below have compiled proofs without added axioms or admitted goals. The table lists one principal declaration per result; the full lists, and notes on how each Lean statement matches the text, are in \path{papers/FORMALIZATION.md}, which together with \path{papers/DIAGNOSTICS.md} is in the later revision~\cite{docs}.
\begin{center}
\begin{tabular}{@{}p{0.36\linewidth}p{0.59\linewidth}@{}}
Result & Lean declaration\\[3pt]

Theorem~\ref{thm:criterion}: geometric criterion & \path{Szkraw.geometric_criterion}\\[3pt]
Lemma~\ref{lem:parameters}: coefficient parametrization & \path{Szkraw.geometric_parameters}, \dots\\[3pt]
Lemma~\ref{lem:basis}: basis and positivity & \path{Szkraw.squarefreeBasis}, \dots\\[3pt]
Lemma~\ref{lem:trace}: positive associative inner product & \path{Szkraw.squarefree_trace_diagonal}, \dots\\[3pt]
Lemma~\ref{lem:three}: three-term multiplication & \path{Szkraw.elementary_three_term}, \dots\\[3pt]
Proposition~\ref{prop:subalgebra}: positive subalgebra & \path{Szkraw.positive_subalgebra_full}\\[3pt]
Lemma~\ref{lem:normalization}: positive normalization & \path{Szkraw.towerCharacter}, \dots\\[3pt]
Lemma~\ref{lem:finite}: finite realization and triangle support & \path{Szkraw.finite_realization}, \dots\\[3pt]
Lemma~\ref{lem:margin}: uniform positive margin & \path{Szkraw.kraw_triangle_margin}\\[3pt]
Theorem~\ref{thm:perturbation}: perturbation & \path{Szkraw.perturbation_criterion}\\[3pt]
Corollary~\ref{cor:nonclassical}: systems outside the geometric family & \path{Szkraw.nongeometric_perturbation_example}, \dots\\[3pt]
Corollary~\ref{cor:kraw}: full positive Krawtchouk range & \path{Szkraw.full_krawtchouk_binomial_range}, \path{Szkraw.krawtchouk_range_sharp}, \dots\\[3pt]
Corollary~\ref{cor:additional}: dual $q$-Krawtchouk case & \path{Szkraw.additional_geometric_class}\\[3pt]
Proposition~\ref{prop:source}: Koelink's relation (3.2) and the recurrence of his Theorem~7.1, transcribed by hand; agreement checked in Lean & \path{Szkraw.koelink_relation_32}, \path{Szkraw.koelink_recurrence_71}\\[3pt]
Theorem~\ref{thm:szwarc}: Szwarc's quadratic system & \path{Szkraw.szwarc_full}, \dots\\[3pt]
Lemma~\ref{lem:edge}: edge idempotents & \path{Szkraw.edge_idempotents}, \dots\\[3pt]
Lemma~\ref{lem:qthree}: three-term action & \path{Szkraw.three_term_action}, \dots\\[3pt]
Lemma~\ref{lem:qmargin}: margin for the quadratic system & \path{Szkraw.szwarc_margin}\\[3pt]
Lemma~\ref{lem:stability}: stability of a reference with a margin & \path{Szkraw.stability}\\[3pt]
Theorem~\ref{thm:symmetric}: symmetric criterion & \path{Szkraw.symmetric_criterion}\\[3pt]
Corollary~\ref{cor:symmetric-new}: symmetric systems outside both families & \path{Szkraw.symmetric_new}, \path{Szkraw.symmetric_open}\\[3pt]
Theorem~\ref{thm:low}: low products & \path{Szkraw.low_products}\\[3pt]
Corollary~\ref{cor:convex-low}: convex profiles & \path{Szkraw.convex_low}, \dots\\[3pt]
Theorem~\ref{thm:all}: all products & \path{Szkraw.all_products}\\[3pt]
Corollary~\ref{cor:heavy}: heavy top step & \path{Szkraw.heavy_top}\\[3pt]
Lemma~\ref{lem:reversal}: reversal identity & \path{Szkraw.reversal_identity}\\[3pt]
Theorem~\ref{thm:deep}: deep reduction & \path{Szkraw.deep_formula}, \dots\\[3pt]
Theorem~\ref{thm:palindromic}: palindromic criterion & \path{Szkraw.palindromic}\\[3pt]
Corollary~\ref{cor:palin}: monotone palindromic systems & \path{Szkraw.palin_mono}, \dots\\[3pt]
Lemma~\ref{lem:corner-value}: corner value & \path{Szkraw.corner_value}\\[3pt]
Lemma~\ref{lem:corner}: corner identity & \path{Szkraw.corner_identity}\\[3pt]
Theorem~\ref{thm:uniform}: forced equalities for uniform criteria & \path{Szkraw.forced_equalities}, \dots\\[3pt]
Lemma~\ref{lem:nr}: N\"orlund--Rice bound & \path{Szkraw.norlund_rice}\\[3pt]
Corollary~\ref{cor:analytic-rules}: analytic rules & \path{Szkraw.analytic_forced_equalities}, \dots\\[3pt]
Proposition~\ref{prop:families}: families depending on $N$ & \path{Szkraw.families_sign}, \dots\\[3pt]
Lemma~\ref{lem:topmass}: lowering the top mass & \path{Szkraw.lower_measure}, \dots\\[3pt]
Proposition~\ref{prop:interior}: interior points near the symmetric endpoint & \path{Szkraw.interior_measure}, \dots\\[3pt]
Proposition~\ref{prop:thin}: thinness near the symmetric endpoint & \path{Szkraw.thinness}\\[3pt]
Corollary~\ref{cor:radius-order}: optimal order of the stability radius & \path{Szkraw.radius_order}\\[3pt]
Section~\ref{sec:limits}: terminal-row failures at fixed $r$ & \path{Szkraw.terminal_fails}\\[3pt]
Section~\ref{sec:limits}: the rule $\beta=1$, $\gamma=0$ at $r=4/5$ and $r=19/20$ & \path{Szkraw.rule_fails_4_5}, \path{Szkraw.rule_fails_19_20}\\[3pt]
Section~\ref{sec:limits}: first-order derivatives along the rule $\beta=1$, $\gamma=0$ at $r=4/5$ & \path{Szkraw.PolyCol.rule_deriv_4_5}\\[3pt]
Section~\ref{sec:limits}: the one-kink counterexamples with $N\leq100$ and the exact log-concave-increment counterexamples & \path{Szkraw.kink20_2_5_fails}, \path{Szkraw.kink40_1_3_fails}, \path{Szkraw.kink40_3_10_fails}, \path{Szkraw.kink60_1_4_fails}, \path{Szkraw.kink80_2_7_fails}, \path{Szkraw.kink100_1_5_fails}, \path{Szkraw.stepInc_logconcave}, \path{Szkraw.logconcave32_fails}, \path{Szkraw.logconcave80_fails}, \path{Szkraw.logconcave90_fails}\\[3pt]
Section~\ref{sec:open}: the convex-profile and log-convex-increment counterexamples & \path{Szkraw.kink14_fails}, \path{Szkraw.kink20_2_5_fails}, \path{Szkraw.cvxInc_logconvex}, \path{Szkraw.logconvex18_fails}\\[3pt]
Introduction: the parameter convention of~\cite{szwarc} & \path{Szkraw.szwarc_footnote}\\[3pt]
Remarks after Corollary~\ref{cor:symmetric-new}: $K_{2,1}$ perturbed & \path{Szkraw.k21_perturbation}\\[3pt]
Remarks after Corollary~\ref{cor:heavy} & \path{Szkraw.heavy_top_sharp}, \path{Szkraw.heavy_top_scope}, \path{Szkraw.convex_superadditive}\\[3pt]
End of Section~\ref{sec:low}: reversal obstructions and Miclo's condition & \path{Szkraw.g111}, \path{Szkraw.g222}, \path{Szkraw.reversal_obstruction_kraw}, \path{Szkraw.reversal_obstruction_szwarc}, \path{Szkraw.miclo_condition_fails}\\[3pt]
After Theorem~\ref{thm:uniform}: known directions, independence, open sets of rules & \path{Szkraw.known_directions_satisfy}, \path{Szkraw.no_open_set_of_rules}\\[3pt]
Section~\ref{sec:limits}: polynomial dependence along the terminal row & \path{Szkraw.terminal_polynomial}\\[3pt]
Section~\ref{sec:limits}: the class of convex profiles with log-concave increments & \path{Szkraw.logconcave_increment_class}\\[3pt]
Section~\ref{sec:open}: sign conditions on the shape & \path{Szkraw.sign_pattern}, \path{Szkraw.kraw_mem_signClass}, \path{Szkraw.signClass_zero_zero}
\end{tabular}
\end{center}
Not formalized, or formalized only in part: for Proposition~\ref{prop:source}, Koelink's relation (3.2) and the recurrence of his Theorem~7.1 were transcribed by hand into Lean, and Lean checks that they agree with \eqref{eq:algebra} and with the recurrence of \eqref{eq:parameters}; it does not check that the transcription is faithful to~\cite{koelink}. His type-B Hecke-algebra construction, his spherical functions and the dual $q$-Krawtchouk interpretation are not formalized; no other result depends on this identification. The comparisons with Szwarc's infinite-system conditions and with Miclo's conditions in Sections~\ref{sec:low} and~\ref{sec:open} are not formalized, except the failure of Miclo's monotonicity condition for the symmetric Krawtchouk chain with $N>2$ (\path{Szkraw.miclo_condition_fails}, with the condition transcribed as stated here). The subdivided-cube interpretation and the $O(N)$ and $O(N^2)$ operation counts are not formalized. In Section~\ref{sec:limits}, the following are not formalized: the reduction of arbitrary finite systems to the chosen normalization (the Lean development starts from the normalized recurrence); the equispaced-interpolation interpretation of $\Omega$; the exact values $\Lambda_N=4,10,\ldots,165$ and the affine-normalized top-mass direction with its non-uniform-rule observations (the mapped top-mass results give a sufficient bound and general fixed-coordinate derivatives, not these computed maxima or the displayed affine-normalized formulas; explicit binomial spectral data are formalized separately in \path{Szkraw.BinomialMeasure}); the one-kink example with $N=500$, which is checked in exact arithmetic by \path{smallr_probe.py}, but whose documented Lean-kernel evaluation did not finish within 50 minutes; and the exact contour representation of the first-order change along the terminal row, with the accompanying saddle-point discussion, which is heuristic. The 80-/150-digit examples in Section~\ref{sec:limits}, the 150-/300-digit example and the unsuccessful sampled search in Section~\ref{sec:open} are numerical evidence, not exact sign certificates.

\begin{thebibliography}{9}

\bibitem{askey} R. Askey, \emph{Three notes on orthogonal polynomials}, Stichting Mathematisch Centrum, Afdeling Toegepaste Wiskunde, report TW~114/69, Amsterdam, October 1969; Part~II: \emph{Linearization of the product of orthogonal polynomials II}, \url{https://ir.cwi.nl/pub/7741}.
\bibitem{eagleson} G. K. Eagleson, \emph{A characterization theorem for positive definite sequences on the Krawtchouk polynomials}, Australian Journal of Statistics \textbf{11} (1969), no.~1, 29--38, \href{https://doi.org/10.1111/j.1467-842X.1969.tb00004.x}{\nolinkurl{doi:10.1111/j.1467-842X.1969.tb00004.x}}.
\bibitem{koelink} H. T. Koelink, \emph{$q$-Krawtchouk polynomials as spherical functions on the Hecke algebra of type B}, Transactions of the American Mathematical Society \textbf{352} (2000), no.~10, 4789--4813, \href{https://doi.org/10.1090/S0002-9947-00-02588-5}{\nolinkurl{doi:10.1090/S0002-9947-00-02588-5}}; preprint arXiv:q-alg/9605040v1 (1996), \url{https://arxiv.org/abs/q-alg/9605040v1}. Theorem and equation locators refer to the published version; the theorem numbers agree with the preprint, where relation (3.2) is numbered (3.3).
\bibitem{miclo} L. Miclo, \emph{On the hypergroup property}, Annales de la Facult\'e des Sciences de Toulouse, s\'erie~6, \textbf{26} (2017), no.~2, 417--435, \href{https://doi.org/10.5802/afst.1539}{\nolinkurl{doi:10.5802/afst.1539}}.
\bibitem{miclo2015} L. Miclo, \emph{On the Markov commutator}, preprint arXiv:1505.02540 (2015), \url{https://arxiv.org/abs/1505.02540}.
\bibitem{szwarc} R. Szwarc, \emph{Orthogonal polynomials and Banach algebras}, in \emph{Inzell Lectures on Orthogonal Polynomials}, W. zu Castell, F. Filbir, and B. Forster (eds.), Advances in the Theory of Special Functions and Orthogonal Polynomials, vol.~2, Nova Science Publishers, 2005, 103--139; \S5, Problem~2. Author version (2004, 97--133; Problem~2 on printed p.~125): \url{https://www.math.uni.wroc.pl/~szwarc/pdf/inzell.pdf}.
\bibitem{szwarc2003} R. Szwarc, \emph{A necessary and sufficient condition for nonnegative product linearization of orthogonal polynomials}, Constructive Approximation \textbf{19} (2003), 565--573, \href{https://doi.org/10.1007/s00365-002-0524-z}{\nolinkurl{doi:10.1007/s00365-002-0524-z}}.
\bibitem{P1} L. P. Polzer, \emph{A log-concavity criterion for finite nonnegative linearization}, companion paper (2026).
\bibitem{P2} L. P. Polzer, \emph{Increasing ratios in finite nonnegative linearization}, companion paper (2026).
\bibitem{code} L. P. Polzer, \emph{Companion code for finite nonnegative linearization} (2026), revision \texttt{a25bf9a} (full identifier in the URL), \url{https://github.com/skypher/szkraw/tree/a25bf9a148047a6d909ae752e1ed1f8452aea394}.
\bibitem{docs} L. P. Polzer, \emph{Companion code for finite nonnegative linearization}, revision \texttt{b4c40c5} with the split documentation \path{papers/FORMALIZATION.md}, \path{papers/DIAGNOSTICS.md} and the fresh-clone record \path{runs/verify/RESULTS.md} (full identifier in the URL), \url{https://github.com/skypher/szkraw/tree/b4c40c54ec479a1a2e3df2739df35124f36f4673}.
\end{thebibliography}
\end{document}
